\documentclass[11pt]{article}
\usepackage[margin=1.15in]{geometry}
\usepackage{amsmath,amssymb,amsthm,mathtools}
\usepackage{xcolor}
\usepackage[colorlinks=true,linkcolor=blue!60!black,citecolor=blue!60!black,urlcolor=blue!60!black]{hyperref}

\newtheorem{theorem}{Theorem}[section]
\newtheorem{proposition}[theorem]{Proposition}
\newtheorem{lemma}[theorem]{Lemma}
\newtheorem{corollary}[theorem]{Corollary}
\newtheorem{cproposition}[theorem]{Computational Proposition}
\theoremstyle{definition}
\newtheorem{definition}[theorem]{Definition}
\theoremstyle{remark}
\newtheorem{remark}[theorem]{Remark}

\newcommand{\C}{\mathbb{C}}
\newcommand{\R}{\mathbb{R}}
\newcommand{\dg}{\dagger}
\newcommand{\dd}{\ddagger}
\newcommand{\ad}{\operatorname{ad}}
\newcommand{\ob}{\mathfrak o}
\newcommand{\Tr}{\operatorname{Tr}}
\newcommand{\Res}{\operatorname*{Res}}
\newcommand{\AK}{\mathcal A_K}
\newcommand{\AKred}{\mathcal A_K^{\rm red}}
\newcommand{\Ap}{\mathcal A_+}
\newcommand{\PiE}{\Pi_{\ker\ad_{h_0}}}
\newcommand{\Born}{\mathcal B}

\title{Interaction Obstructions in Einstein--Weyl Gravity:\\
Cubic Protection, Second-Order Metric Failure, and Krein Visibility}
\author{GPT-5.6.sol\thanks{OpenAI model and principal programme author.}
\and Asger Alstrup Palm\thanks{Programme orchestrator and corresponding human contact; Honorary Professor, DTU Compute, Technical University of Denmark. \texttt{asger@area9.dk}.}}
\date{Open-repository preprint, August 2026}

\begin{document}
\maketitle

\begin{abstract}
We study whether the two Lorentz-covariant free completions of
scalar-free Einstein--Weyl gravity extend to the interacting theory.  The
positive pseudo-Hermitian completion uniformly quarter-turns the complete
massive spin-$2$ multiplet, whereas the conventional Krein completion
retains standard gravitational reality and assigns negative signature to
that multiplet.  We first prove a cubic-protection theorem.  The exact
Einstein consistent truncation implies that every tree amplitude with
exactly one massive external field and otherwise massless gravitons
vanishes, while the equal-mass $MMM$ vertex has no real
$1\leftrightarrow2$ shell.  Hence the first-order positive-metric
obstruction is zero.

At second order the physical process $MM\to Mh$ is open.  In an
independent second-order Einstein-frame formulation we compute nonzero
on-shell $MMM$ and $MMh$ amplitudes and the exact pole-normalized massive
residue $\sqrt3/8$, proving that the four-point amplitude is not
identically zero.  A complete calculation in the original fourth-order
theory --- including the quartic contact term, all $s,t,u$ exchanges, the
physical adjoint, and gauge and equation-of-motion checks --- gives at an
interior rational physical point an exact nonzero rational certificate.
This complete physical branch-changing amplitude is the central new
input; the surrounding auxiliary-field and pseudo-Hermitian formalisms
are established tools.
We prove an operator identity equating the anti-Hermitian part of the
stationary second-order Born operator to the full second-order
metric-deformation source on the free-energy shell.  It follows that
\begin{equation*}
\PiE\left[(v_2^{\dg}-v_2)
+\frac12[R_1,v_1+v_1^{\dg}]\right]\ne0 .
\end{equation*}
The canonical analytic deformation based at the Paper-IV positive real
form therefore fails at second order for the complete tree-level vertex
content through quartic order.

On physical BRST cohomology, canonical second quantization of the free
Krein symmetry gives the particle-number grading $J_F=(-1)^{N_M}$.
The gravitational obstruction survives cohomology and is not null:
\begin{equation*}
\mathcal O^\sharp=\mathcal O,
\qquad
\Tr(\mathcal O^\sharp\mathcal O)=-8\AK^2\ne0 .
\end{equation*}
Unlike the perfect-square scalar theory, Einstein--Weyl gravity has no
nontrivial uniform abelian charge grading capable of placing the
obstruction in a one-sided null ideal.  The result excludes both an
analytic continuation of the canonical positive completion and the
conventional massive-number grading as a null-sector probability
mechanism.  It does not exclude indefinite Krein evolution itself, nor
nonlocal, non-factorizing, nonperturbative, or conformal-boundary
completions.  An exact scalar \texttt{REDUCED-MODE} calculation also
shows that the ordinary axis-compatible independent-mass regulator does not
cancel its complete final-state collinear logarithmic response.  That
is not a gravitational theorem and changes none of the tree-level results
here; it sharpens the extra infrared/asymptotic-state input any
Bateman--Turok-type gravitational completion would have to supply.
A finite coherent-projector calculation cancels that scalar coefficient while
preserving idempotence, but it does not derive the transport dynamically and
does not transfer to the gravitational BRST carrier.  An exact
scalar countermodel shows that CCR preservation, projector idempotence, trace
preservation, and endpoint exchange symmetry still leave three endpoint
parameters.  The scalar $1/48$ value is compatible but not selected by those
identities, sharpening rather than removing the need for an independently
constructed continuum transport.  The exact full off-resonant scalar
composition reduces the flat three-parameter endpoint freedom to one soft
normalization on its declared three-dimensional chart, but its measured
cross-Gram remains $-(1/2)\,dr/r$ and hence logarithmically non-trace-class.
The full normalization calculation fixes the unprojected scalar response to
$1/48$ per pair, but its logarithmic rows pair fixed-vacuum charges $+1$ and
$-1$; the published one-sided projection instead makes that residue zero.
The exact scalar vacuum-orbit operator makes those rows neutral and preserves
$1/48$, but simultaneously neutralizes the squeeze that the fixed-vacuum
negative grading discarded; the charge derivation does not descend after
setting the orbit operator to one.  The missing full zero-mode trace therefore
still decides the physical answer and supplies no gravitational or causal
normalization.
The same Appendix-C scalar squeeze also fails to define a vector in the
ordinary massless positive Fock--Krein topology: its two-particle norm density
diverges as the inverse infrared cutoff, and its pair block is not
Hilbert--Schmidt.  This remains a scalar
\texttt{LOCAL-ALGEBRAIC}, \texttt{REDUCED-MODE} carrier obstruction.  It
neither transfers to the metric BV complex nor excludes an explicitly
extended/non-Fock scalar representation.
The full scalar exponential admits an explicit infrared-weighted vacuum
carrier, but its inverse metric is unbounded, its thermodynamic representation
is inequivalent.  Its covariant squeeze does admit a cross-Krein core
implementation and cyclic finite-rank trace, but a finite normalized orbit
trace kills localized projectors and the transported vacuum projector has
exponentially growing trace norm.  No normal thermodynamic Born trace follows.
A semifinite scalar construction nevertheless gives normalized finite-detector
weights under the weak ghost hypothesis by conditioning on a finite incoming
projection; its normalized corner functional is noncyclic and supplies neither
a thermodynamic state nor Eq.~(19).
On the available zero-mode-completed logarithmic sector, an exact scalar
detector pushforward is finite rank at every fixed soft cutoff and remains in
that semifinite ideal after the weighted squeeze.  Projector idempotence forces
opposite hard and soft box traces, but their cancellation does not control the
positive trace ideal: for $N$ logarithmic cells the order-$\lambda$ coefficient
has squared trace norm $N/12$ and the order-$\lambda^2$ coefficient has trace
norm $N/24$.  Thus this sector has no uniform trace-class soft limit without an
additional full-pushforward cancellation, hard matching operator, or local
non-normal renormalized weight.  The result neither supplies the missing
scalar Eq.~(19) sectors nor transfers to the gravitational BV carrier.
The complete scalar signed quadratic pushforward provides that cancellation
at order $\lambda$: the off-resonant symplectic phase cancels the oscillatory
inverse phase, every endpoint-dangerous row vanishes, and the parent cross-Gram
is zero.  Its mixed-sign rows satisfy the cross-CCR Ward identity, while the
unique zero-mode dressing makes the finite-mode generator neutral.  Hence the
finite-mode quadratic carrier obeys Eq.~(19) through order $\lambda$ with
$Q_1=0$.  The conditional $1/48$ is absent from the completed public quadratic
map.  This does not determine the physical scalar response because the
S-matrix transition block is a distinct operator.  The squeezed-vacuum and
dynamical-zero-mode projector trace, continuum domains, and higher orders
remain open within Eq.~(19).
This remains solely a scalar \texttt{LOCAL-ALGEBRAIC},
\texttt{REDUCED-MODE} result.
The covariant Appendix-C squeeze conjugates the finite-core scalar detector
and Born operator together, so finite-rank cyclicity preserves this zero; bare
squeezed-vacuum pair occupation is not an additive Eq.~(19) contribution.
The missing thermodynamic trace domain still limits Eq.~(19), while the
operator distinction independently forbids reading its zero as an S-matrix
coefficient.
Finite paired scalar processes nevertheless have a spectator-stable
non-normal cylinder functional, on which the zero coefficient survives; this
is not an inclusive LSZ or gravitational construction.
For a scalar orthogonal inclusive detector, higher composite amplitudes do not
enter the leading transition weight, but the input kernel must still be the
physical S-matrix block.  The complete signed zero belongs instead to the
distinct scalar Eq.~(19) pushforward.  The scalar five-point S-matrix response
is $+1/512$ per pair and $+3/512$ in total, or $1/48$ per pair after Born
normalization.  The axis-compatible physical virtual response is zero, leaving
$+3/512$ uncancelled.  The public $R_t$ response is also zero but is a distinct
Eq.~(19) object, not a physical summand.  An independently normalized
hard/dressed response $-3/512$ is still missing.  This scalar object ledger
neither completes its NLO probability nor transfers a coefficient or state
construction to gravity.
On a regulated positive scalar quotient, physical-shell pseudo-unitarity
forces that coefficient: the real-column norm $1/16$ implies the hard survival
response $-1/16$, hence absolute $-3/512$ after Born normalization.  An exact
finite isometric witness exists, but the statement is conditional on a
continuum dressed scalar S-matrix and does not transfer to gravity.
On the ordinary logarithmic scalar carrier that strong Møller limit is
obstructed by fixed-norm shells escaping to the collinear endpoint.  A
regulator-pulled-back abstract endpoint fibre retains the leading-log
isometry, but has no certified local or gravitational affiliation.
\end{abstract}

\section*{AI authorship and computational accountability}
GPT-5.6.sol developed the derivations, computational checks,
claim-boundary analysis, and manuscript.  Asger Alstrup Palm commissioned
and orchestrated the programme
and is the corresponding human contact.  The project is an experiment in
whether a non-specialist human, using AI systems as research collaborators,
can organize falsifiable computer-assisted mathematical physics.  The exact
status of each repository check is stated where it is used, and no check is
promoted beyond its documented scope.  The manuscript remains subject to
expert review.

\tableofcontents

\section{Introduction}\label{sec:intro}

Scalar-free Einstein--Weyl gravity is the simplest four-dimensional
quadratic-curvature theory whose flat-space physical spectrum consists of
the ordinary massless graviton and one additional massive spin-$2$
multiplet.  In the curvature convention used here its action is
\begin{equation}\label{eq:action}
S[g]=\int d^4x\,\sqrt{-g}\left(
c_1R+\alpha R_{\mu\nu}R^{\mu\nu}+\beta R^2\right),
\qquad \alpha=-3\beta,
\end{equation}
with $M^2=c_1/\alpha>0$ in the split phase.  Up to the four-dimensional
Euler density, the quadratic term is Weyl-squared.  The extra spin-$2$
branch has the opposite kinetic sign, the familiar higher-derivative
ghost of quadratic gravity \cite{Stelle,Salvio}.  Neither this spectrum
nor its nonlinear auxiliary-field reformulation is new: quadratic
gravity was put in a local second-order Einstein--scalar--massive-tensor
form by Hindawi, Ovrut, and Waldram \cite{HindawiOvrutWaldram}, and the
scalar-free massive-tensor sector was developed in detail by Magnano and
Soko\l{}owski \cite{MagnanoSokolowski}.

Paper~IV of this series \cite{paperIV} did not choose between proposed
interacting cures.  It classified the \emph{free} covariant completions.
After symplectic reduction, all five massive polarizations form one
irreducible massive spin-$2$ representation.  Covariance therefore
allows no helicity-by-helicity compromise.  There are two real-form
classes in the stated free, quasifree, mode-local setting:
\begin{enumerate}
\item a positive pseudo-Hermitian completion in which the entire massive
multiplet has quarter-turned reality; and
\item a Krein completion with standard gravitational reality, positive
massless signature, and uniform negative massive signature.
\end{enumerate}
They determine the same free complex spectral functional but different
adjoints and probability structures.  The natural interacting question
is therefore sharp:
\begin{quote}
Does either covariant free completion survive the tree interaction
through quartic order?
\end{quote}

Paper~V \cite{paperV} supplied the deformation complex used here to ask
the positive part of this question.  Perturbative pseudo-Hermitian
metrics, their homogeneous ambiguities, and interacting field-theory
constructions have a substantial prior literature
\cite{MostafazadehMetric,BBJQFT}; we do not claim the order-by-order
formalism itself as new.  If $\rho_0$ is the free Dyson map and
$v_n$ are the transported interaction vertices, an analytic correction
$\rho_\zeta=\exp[-\tfrac12\sum_n\zeta^nR_n]\rho_0$ must solve
\begin{align}\label{eq:defcomplex}
[h_0,R_1]&=v_1^{\dg}-v_1,\nonumber\\
[h_0,R_2]&=(v_2^{\dg}-v_2)
+\frac12[R_1,v_1+v_1^{\dg}].
\end{align}
On a degenerate free-energy shell the right-hand side must have zero
projection onto $\ker\ad_{h_0}$.  That shell projection is the metric
obstruction.  In the perfect-square scalar, first-order protection was
followed by a second-order branch-changing obstruction; at the exact
massless $O(1,1)$ boundary, however, the offending component could become
one-sided in continuous charge and null in the generalized Born trace.

\paragraph{Prior amplitudes and the precise novelty claim.}
The geometric premise behind the cubic selection rule is also known:
Einstein metrics form an exact Bach-flat subsector.  Related amplitude
work proves that Ricci-curvature deformations leave all-Einstein-graviton
tree amplitudes unchanged and that the corresponding pure conformal
amplitudes vanish \cite{DonaEtAl}.  Lemma~\ref{lem:onemassive} below is
the mass-eigenfield, one-massive-leg consequence of that established
consistent truncation, not a new discovery of the Einstein subsector.
Conformal-gravity BRST/LSZ quantization and its indefinite scattering
problem have likewise been studied directly \cite{KuboKuntz}, and
conformal or mass-deformed ghost-state amplitudes are known in adjacent
conformal and supersymmetric theories
\cite{AdamoNakachTseytlin,JohanssonMogullTeng}.  Those works concern the
conformal/Jordan phase or different amplitude sectors, not the regular
split-phase process studied here.

The closest comparison is Kuntz's PT-symmetric quadratic-gravity
proposal \cite{KuntzPTGravity}.  It uses the same scalar-free quadratic
action and, up to convention, the same complex rotation of the complete
massive spin-$2$ field.  It constructs the free BRST physical subspace
and solves an interacting pseudo-Hermitian metric in a mode-truncated
Hamiltonian in which spin and momentum labels and the exact gravitational
vertex coefficients are suppressed.  It does not evaluate a complete
degenerate physical scattering shell.  The calculation below supplies
precisely that missing stress test.

To our knowledge, the narrow new input of this paper is the first
complete exact calculation of the physical mass-eigenfield process
$MM\to Mh$ in scalar-free Einstein--Weyl gravity, together with its
independent complete massive-pole residue, and the use of that physical
block to produce an explicit metric-deformation cocycle.  The novelty is
the verified gravitational source, not the auxiliary formulation, the
general deformation equations, the elementary external-leg phase count,
or indefinite-metric algebra.  Propositions~\ref{prop:bornidentity} and
\ref{prop:covariantborn} make the amplitude-to-cocycle step an operator
identity rather than a phase-based inference.

Gravity has the same perturbative-order pattern and a different
mechanism:
\begin{equation}\label{eq:pattern}
\boxed{\ob_+^{(1)}=0,\qquad \ob_+^{(2)}\ne0.}
\end{equation}
The cubic protection follows from the exact Einstein subsector together
with equal-mass kinematics.  The second-order obstruction is carried by
the physically open process
\begin{equation}\label{eq:process}
M+M\longrightarrow M+h,
\end{equation}
which changes massive particle number by one.  The same process also
separates gravity from the scalar boundary mechanism.  Covariance and
the free signature fix the one-particle symmetry; its natural
particle-number-diagonal, cluster-multiplicative lift is $(-1)^{N_M}$.
But $\mathbb Z_2$-odd does not imply null:
the product of two odd operators is neutral and its quadratic trace is
nonzero.

\paragraph{Results.}
The proof has four independent layers.
\begin{enumerate}
\item Section~\ref{sec:phase} proves a diagram-level phase theorem: a
connected amplitude with $E_M$ external massive legs changes by
$(-i)^{E_M}$ under the uniform quarter-turn, including the compensating
sign of every internal massive propagator.
\item Section~\ref{sec:cubic} proves cubic protection.  The Einstein
consistent truncation kills $Mhh$, while $MMM$ has no real physical
$1\leftrightarrow2$ shell.
\item Sections~\ref{sec:factorization}--\ref{sec:full} prove the
second-order obstruction twice: first by an independent nonzero
Einstein-frame factorization residue, then by a complete exact physical
four-point calculation in the original fourth-order theory, joined to
the deformation complex by Proposition~\ref{prop:bornidentity}.
\item Sections~\ref{sec:grading}--\ref{sec:brst} prove conventional
Krein visibility: the canonical multiplicative Fock lift is fixed in the
stated class, the exact block is not null, no scalar-like abelian charge
exists, and the block is not BRST-exact on physical cohomology.
\end{enumerate}

\paragraph{Comparison with the scalar obstruction.}
The parallel and the difference are summarized by
\begin{center}
\begin{tabular}{p{0.42\linewidth}|p{0.42\linewidth}}
\hline
Perfect-square scalar & Einstein--Weyl gravity\\
\hline
K\"all\'en zero and reality-factor protection & Einstein consistent
truncation and equal-mass kinematics\\
$H+L\to L+L$ & $M+M\to M+h$\\
$O(1,1)$ one-sided null ideal at the exact boundary & no nontrivial
uniform abelian charge grading\\
obstruction can become probability-invisible & obstruction remains
non-null for the conventional covariant Krein form\\
\hline
\end{tabular}
\end{center}

\paragraph{Strength of the claims.}
``Metric failure'' means failure of the analytic order-by-order
deformation of the canonical covariant positive free metric in the
operator class of \eqref{eq:defcomplex}.  It is not an absolute theorem
against every positive metric, and no complex interacting spectrum is
proved.  ``Krein visibility'' has the precise finite-block meaning
$\mathcal O^\sharp\mathcal O\ne0$ and
$\Tr(\mathcal O^\sharp\mathcal O)\ne0$ for the canonical
particle-number-diagonal, cluster-multiplicative grading.  It is a
non-nullity statement,
not a claim that an indefinite trace is a positive probability.  In
particular, the paper does not rule out an indefinite Krein dynamics;
it rules out the naive parity/null-sector mechanism as a replacement for
a positive interacting probability calculus.

\paragraph{Verification.}
The labels G13--G18 denote entries in the machine-verification table of
Appendix~\ref{app:checks}; they are not additional assumptions or
mathematical objects.  The checks use exact SymPy or rational arithmetic.
The exhaustive $5^3\times2$ polarization scan once proposed as G16 is
not used:
the factorization residue, exact real-shell certificate, crossing, Ward,
and gauge checks already close the proof.  The scan remains optional
regression hardening.

\section{Free input, conventions, and the deformation class}
\label{sec:free}

\subsection{Physical one-particle spectrum}

We use signature $(+,-,-,-)$ in the main text.  Around flat space the
reduced physical one-particle space is
\begin{equation}\label{eq:oneparticle}
\mathcal H_{\rm phys}^{(1)}=
\mathcal H_{h,+2}\oplus\mathcal H_{h,-2}\oplus\mathcal H_M,
\qquad \dim\mathcal H_M=5.
\end{equation}
The massless summands carry helicities $\pm2$ on $p^2=0$; the massive
summand carries the spin-$2$ representation on $p^2=M^2$.  Paper~IV
derives this by explicit diffeomorphism reduction rather than by a
covariant propagator count.  Its result is a healthy massless sector and
a uniformly ghost-signed massive sector.

\subsection{Two free real forms}

In the positive class the original massive mass eigenfield is
$\dd$-anti-Hermitian.  We write the positive-frame observable as
\begin{equation}\label{eq:quarterturn}
\widehat M_{\mu\nu}=iM_{\mu\nu},
\qquad M_{\mu\nu}=-i\widehat M_{\mu\nu}.
\end{equation}
The same quarter-turn acts on every one of the five polarizations.  The
massless graviton retains standard reality.  In the conventional Krein
class both fields retain standard gravitational reality, while the
one-particle fundamental symmetry has signature
\begin{equation}\label{eq:freesignature}
J^{(1)}_{\rm free}=I_2\oplus(-I_5),
\qquad \text{signature }(+,+;-,-,-,-,-).
\end{equation}

The uniformity is not optional.  Since $\mathcal H_M$ is an irreducible
massive spin-$2$ representation, Schur's lemma makes every covariant
Hermitian form on its fiber a scalar multiple of the identity.  Rotating
only selected helicities would violate Lorentz covariance
\cite{paperIV}.

\subsection{Regulated physical Fock space and shell projection}

As in Paper~V, the cohomological equation is first interpreted in a
periodic spatial box.  We retain finitely many nonzero momentum modes,
impose an ultraviolet and a finite-particle cutoff, and omit the
massless zero mode.  Normal ordering removes vacuum contractions, and
we work with connected transition kernels.  On this regulated physical
Fock space $h_0$ has point spectrum with finite-rank eigenspaces.  If
$P_E$ projects onto free energy $E$, then
\begin{equation}\label{eq:kernelproj}
\PiE S=\sum_E P_ESP_E,
\end{equation}
and $[h_0,R]=S$ has a solution only if this projection vanishes.

The rational point used below lies on the momentum lattice whenever
$ML/(2\pi)$ is a multiple of $100$.  The ultraviolet cutoff may be
chosen above every external and internal momentum occurring in the
calculation.  We take a sequence of such boxes after establishing the
finite-volume obstruction.  Omitting the zero mode is the infrared
prescription; the exhibited momenta remain uniformly nonsoft along this
sequence.  Thus neither the matrix element nor its internal denominators
depends on a soft limiting prescription.

\subsection{Perturbative order and the stationary Born operator}

Introduce a formal coupling $\zeta$ by assigning the cubic expansion of
the Einstein--Weyl action order $\zeta$ and its quartic expansion order
$\zeta^2$.  After the Paper-IV free Dyson map, write
\begin{equation}\label{eq:positiveexpansion}
H_+(\zeta)=h_0+\zeta v_1+\zeta^2v_2+O(\zeta^3).
\end{equation}
The second-order metric-deformation source is
\begin{equation}\label{eq:S2definition}
S_2=(v_2^{\dg}-v_2)
+\frac12[R_1,v_1+v_1^{\dg}],
\qquad [h_0,R_1]=v_1^{\dg}-v_1.
\end{equation}
This is the meaning of ``second order'' throughout the paper: one
quartic vertex or two cubic vertices at tree level.

Set $Q_E=1-P_E$ and define the self-adjoint reduced resolvent
\begin{equation}\label{eq:reducedresolvent}
G_E=Q_E(E-h_0)^{-1}Q_E.
\end{equation}
The stationary connected second-order Born operator on the shell is
\begin{equation}\label{eq:Bornoperator}
\Born_{2,+}(E)=P_E\bigl(v_2+v_1G_Ev_1\bigr)P_E.
\end{equation}
At finite volume no $i0$ is required.  In the continuum notation
$G_E$ is the principal-value part of $(E-h_0+i0)^{-1}$.  The possible
$-i\pi\delta(E-h_0)$ contribution is absent here because the regulated
physical cubic shell is empty; Lemma~\ref{lem:fullcubicshell} below
makes this statement on the complete Fock eigenspace, including
spectators.

\begin{proposition}[Born--deformation identity]
\label{prop:bornidentity}
Suppose $P_Ev_1P_E=0$ and choose the canonical off-shell solution of the
first equation in \eqref{eq:S2definition},
\begin{equation}\label{eq:R1spectral}
R_1^{(0)}=\sum_{E_a\ne E_b}
\frac{P_{E_a}(v_1^{\dg}-v_1)P_{E_b}}{E_a-E_b}.
\end{equation}
Then the entire second-order metric source, not only its quartic part,
satisfies
\begin{equation}\label{eq:bornidentity}
\boxed{
P_ES_2P_E
=\Born_{2,+}(E)^{\dg}-\Born_{2,+}(E).}
\end{equation}
Consequently, a nonzero ordinary-adjoint anti-Hermitian Born matrix
element on the shell is a nonzero metric-deformation cocycle.
\end{proposition}

\begin{proof}
Write $D=v_1^{\dg}-v_1$ and $W=v_1+v_1^{\dg}$.  The spectral formula
\eqref{eq:R1spectral} gives
\begin{equation}
P_ER_1^{(0)}Q_E=P_EDG_E,
\qquad
Q_ER_1^{(0)}P_E=-G_EDP_E.
\end{equation}
Since $P_Ev_1P_E=0$, insertion of $Q_E$ between the two cubic vertices
is automatic.  Hence
\begin{align}
\frac12P_E[R_1^{(0)},W]P_E
&=\frac12P_E(DG_EW+WG_ED)P_E\nonumber\\
&=P_E(v_1^{\dg}G_Ev_1^{\dg}-v_1G_Ev_1)P_E.\label{eq:commborn}
\end{align}
Adding $P_E(v_2^{\dg}-v_2)P_E$ gives
\eqref{eq:bornidentity}, because $G_E=G_E^{\dg}$.  Thus the quartic
contact and the two-cubic exchange contribution enter with exactly the
relative combination required by the deformation equation.
\end{proof}

\begin{remark}[Contact, exchange, and the logical bridge]
The two terms in the deformation source have a fixed diagrammatic
meaning.  The shell projection of $v_2^{\dg}-v_2$ is the
anti-Hermitian difference of the canonical quartic vertex, including
the instantaneous terms generated by the Legendre transform.  Equation
\eqref{eq:commborn} is the anti-Hermitian difference of the complete
two-cubic intermediate-state sum.  Their sum, and only their sum, is the
connected second-order Born block.  Thus the four-point amplitude below
is not being substituted for $S_2$ on the basis of its external-leg
phase: Proposition~\ref{prop:covariantborn} first identifies the complete
covariant contact-plus-exchange calculation with \eqref{eq:Bornoperator},
and \eqref{eq:bornidentity} then identifies its anti-Hermitian part with
$P_ES_2P_E$.  The finite-volume principal-value prescription, external
normalizations, identical-particle factors, and removal of disconnected
pieces are stated separately below.
\end{remark}

\subsection{Reduced amplitude, LSZ normalization, and adjoint}

The perturbiner strips the universal Feynman $i$ from the multilinear
action coefficient; denote the result by $\AKred$.  At real momenta and
real linear polarizations it is real.  Let $|i\rangle,|f\rangle$ be
unit-normalized finite-volume asymptotic states.  Matching the covariant
tree expansion to the stationary Born series gives
\begin{equation}\label{eq:LSZmap}
\langle f|\Born_{2,K}(E)|i\rangle
=\AK=\eta_{\rm conv}\,\mathcal N_{fi}\,\AKred,
\qquad \eta_{\rm conv}\in\{\pm1\},\quad \mathcal N_{fi}>0.
\end{equation}
After removing the overall momentum-conservation Kronecker delta,
$\mathcal N_{fi}$ is the product of the nonzero external pole
residues, $(2\omega L^3)^{-1/2}$ factors, polarization normalizations,
and Bose normalizations.  At the displayed point the two incoming
massive momenta are distinct, so their identical-particle factor is one;
the outgoing particles have different species.  The multilinear wave
labels already implement the vertex symmetry factors.  The real sign
$\eta_{\rm conv}$ fixes the convention relating the interaction
Hamiltonian to the Lagrangian Feynman rule.

Equation~\eqref{eq:LSZmap} also explains why stripping the universal
Feynman $i$ does not manufacture the obstruction: the removed factor is
common to forward and reverse $S$-matrix elements, whereas the
anti-Hermitian phase below comes from the three massive quarter-turns.
Real, nonzero LSZ and normalization factors cannot change that phase or
the nonvanishing result.  Exact rational values are quoted for
$\AKred$; operator and trace statements use the normalized coefficient
$\AK$.

\section{The external-massive-leg phase theorem}\label{sec:phase}

The quarter-turn phase must be transported through complete diagrams,
not appended to an already computed external state by slogan.  Internal
massive lines carry the compensating kinetic sign.

\begin{theorem}[External-massive-leg phase]\label{thm:phase}
Let $\mathcal D$ be a connected perturbative diagram written in
mass-eigenfields $(h,M)$, with $E_M$ external massive legs and $I_M$
internal massive lines.  Under \eqref{eq:quarterturn}, its complete
reduced amplitude obeys
\begin{equation}\label{eq:phasetheorem}
\mathcal A_+^{(E_M)}(\mathcal D)
=(-i)^{E_M}\mathcal A_K^{(E_M)}(\mathcal D).
\end{equation}
The statement includes contact terms, exchanges, and the sign of each
massive inverse kernel.
\end{theorem}

\begin{proof}
Let $n_v$ be the number of massive fields incident at vertex $v$.  Field
substitution contributes
\begin{equation}
\prod_v(-i)^{n_v}.
\end{equation}
The massive quadratic form changes sign because
$(-i)^2=-1$; hence its inverse, the internal massive propagator, also
contributes $-1$ per internal massive line.  Every internal massive line
is incident twice, so
\begin{equation}
\sum_v n_v=E_M+2I_M.
\end{equation}
The total factor is therefore
\begin{equation}
(-i)^{E_M+2I_M}(-1)^{I_M}
=(-i)^{E_M}(-1)^{I_M}(-1)^{I_M}
=(-i)^{E_M}.
\end{equation}
Massless propagators are unchanged.  Since the argument is diagramwise,
it applies to their complete sum.  Local invertible changes between
regular mass-eigenfield frames can redistribute contact and exchange
terms but leave the on-shell statement unchanged.
\end{proof}

\begin{corollary}[Phase selection]\label{cor:phase}
For a real standard-frame tree amplitude, even $E_M$ gives a real
positive-frame coefficient and odd $E_M$ gives a purely imaginary one.
Only odd-$E_M$ processes can contribute to the ordinary-adjoint
anti-Hermitian shell source.
\end{corollary}

For the four-point process \eqref{eq:process}, $E_M=3$, and
\begin{equation}\label{eq:mmmhphase}
\mathcal A_+=i\mathcal A_K.
\end{equation}
The sign is tied to the convention $M=-i\widehat M$; reversing the
quarter-turn reverses the displayed sign but not the obstruction's
nonvanishing.

\section{Cubic protection}\label{sec:cubic}

There are two odd-massive cubic structures: $Mhh$ and $MMM$.  The first
is eliminated dynamically, the second kinematically.

\subsection{Einstein consistent truncation}

The scalar-free equations may be written schematically as
\begin{equation}\label{eq:ewfield}
c_1G_{\mu\nu}+\gamma_B B_{\mu\nu}=0,
\end{equation}
where $B_{\mu\nu}$ is the Bach tensor and $\gamma_B$ depends on the
normalization in \eqref{eq:action}.  Every Ricci-flat Einstein metric has
$G_{\mu\nu}=0$ and is Bach-flat.  More generally Einstein metrics are
Bach-flat \cite{LiuLuPope}.  Thus the nonlinear Einstein solution space
is an exact subsector of the Einstein--Weyl equations.

On the regular split branch, the higher-derivative theory admits a local
second-order mass-eigenfield formulation in which the Einstein metric and
the nonlinear massive spin-$2$ field are separate variables
\cite{MagnanoSokolowski}.  Denote them by $(h,M)$ near flat space.  The
field map is algebraic and locally invertible on this branch, and $M=0$
is precisely the Einstein subsector.  These hypotheses are what make the
following LSZ statement stronger than a basis-dependent observation.

\begin{lemma}[One-massive-leg rule]\label{lem:onemassive}
For asymptotic mass eigenfields on the regular split branch,
\begin{equation}\label{eq:onemassive}
\mathcal A_{\rm tree}(M,h,\ldots,h)=0.
\end{equation}
In particular $\mathcal A_3(Mhh)=0$ for every physical polarization.
\end{lemma}

\begin{proof}
Write the massive-field equation in the second-order frame as
\begin{equation}
\mathcal K_M M+J_M^{(2)}[h,h]+J_M^{(3)}[h,h,h]+\cdots=0.
\end{equation}
For every nonlinear vacuum Einstein solution $h_{\rm E}$, the pair
$(h_{\rm E},M=0)$ solves the full theory.  Hence the terms independent of
$M$ vanish on the Einstein graviton shell.  Functional differentiation
and LSZ reduction say that a physical amplitude with one external $M$
and all remaining external legs in the Einstein subsector is zero.
Terms generated by a local invertible change of fields are proportional
to equations of motion or external inverse propagators and vanish under
LSZ, the standard equivalence-theorem mechanism \cite{KOS}.
\end{proof}

The exact perturbiner check evaluates the physical decay point with the
massive particle at rest and two back-to-back massless gravitons.  All
$5\times4=20$ combinations of the five massive polarizations and four
graviton-helicity choices vanish exactly.  This is a check of the
amplitude-level content of Lemma~\ref{lem:onemassive}, not its only
justification.

\subsection{Equal-mass cubic kinematics}

The $MMM$ vertex itself is nonzero at complex on-shell kinematics
(Section~\ref{sec:factorization}).  It nevertheless has no real physical
$1\leftrightarrow2$ shell.  If $p_3=p_1+p_2$ with future-directed
$p_i^2=M^2$, then
\begin{equation}
p_3^2=(p_1+p_2)^2\ge(2M)^2,
\end{equation}
which cannot equal $M^2$.  Equivalently, a particle of mass $M$ cannot
decay into two particles of the same mass.

\subsection{The complete regulated cubic shell}

The Born--deformation identity requires a statement about the complete
degenerate eigenspace, not only the two states used below to display the
four-point obstruction.

\begin{lemma}[Vanishing of the regulated physical cubic shell]
\label{lem:fullcubicshell}
On the regulated physical Fock space of Section~\ref{sec:free}, with the
massless zero mode omitted,
\begin{equation}\label{eq:fullcubicshell}
P_Ev_1P_E=P_Ev_1^{\dg}P_E=0
\end{equation}
for every free energy $E$.  The statement includes disconnected
spectators: the connected cubic vertex cannot act nontrivially inside a
degenerate many-particle shell.
\end{lemma}

\begin{proof}
A normal-ordered cubic matrix element is an active $0\leftrightarrow3$
or $1\leftrightarrow2$ process times Kronecker deltas for particles on
which the vertex does not act.  Positive energy forbids
$0\leftrightarrow3$.  It remains to classify the four physical cubic
field contents.  The $Mhh$ amplitude is zero by
Lemma~\ref{lem:onemassive}; $MMM$ is closed by equal-mass kinematics.
For $MMh$, the channel $h\to MM$ would require
$k^2=(p_1+p_2)^2\ge4M^2$, while $M\to M+h$ would require
$p_M\cdot k=0$ for a future timelike $p_M$ and a future nonzero null
$k$, which is impossible.  The only limiting exception is $k=0$, removed
by the infrared prescription.  Finally, real on-shell $hhh$ kinematics
is collinear, and the physical Einstein three-graviton amplitude
vanishes there.  Spectators do not alter any active momentum or energy
equation, so they multiply these zeros.  Taking the adjoint gives the
second equality in \eqref{eq:fullcubicshell}.
\end{proof}

\begin{theorem}[Cubic protection]\label{thm:cubicprotection}
For scalar-free Einstein--Weyl gravity about flat space, the
Lorentz-covariant positive real form has no first-order physical on-shell
metric obstruction:
\begin{equation}\label{eq:firstzero}
\boxed{\ob_+^{(1)}=0.}
\end{equation}
\end{theorem}

\begin{proof}
By Corollary~\ref{cor:phase}, only odd-massive cubic amplitudes can enter
$\PiE(v_1^{\dg}-v_1)$.  The $Mhh$ amplitude vanishes by
Lemma~\ref{lem:onemassive}; the $MMM$ amplitude has no real physical
shell.  The remaining cubic structures carry an even number of massive
legs.  Equivalently, the stronger Lemma~\ref{lem:fullcubicshell} makes
the complete first-order shell block vanish.
\end{proof}

This protection is structurally different from the K\"all\'en-factor
zero of Paper~V.  It is a consequence of an exact nonlinear consistent
truncation plus massive-particle kinematics.

\section{Independent Einstein-frame factorization rail}
\label{sec:factorization}

The next odd-massive process is four-point $MM\to Mh$.  Before expanding
the original fourth-order action to quartic order, we prove independently
that this amplitude is not identically zero.

\subsection{Second-order action and convention change}

This section follows the exact Einstein-frame formulation of Magnano and
Soko\l{}owski \cite{MagnanoSokolowski}.  Its verification artifact uses
signature $(-,+,+,+)$, only in this section and Appendix~\ref{app:eframe};
thus a massive pole is $P^2=-M^2$.  With $M=1$ and in the notation of that
formulation, the nonlinear massive-field Lagrangian is
\begin{equation}\label{eq:LM}
L_M=\frac12K(Q)+\frac{m^2}{8A}
\left(-\tau^2+\psi_{\mu\nu}\psi^{\mu\nu}
+6A\tau-12A^2\right),
\qquad \psi=g+\Phi .
\end{equation}
Here $Q^\alpha{}_{\mu\nu}=\Gamma^\alpha{}_{\mu\nu}(\psi)
-\Gamma^\alpha{}_{\mu\nu}(g)$,
\begin{equation}
K(Q)=g^{\mu\nu}\left(
Q^\alpha{}_{\mu\nu}Q^\beta{}_{\alpha\beta}
-Q^\alpha{}_{\mu\beta}Q^\beta{}_{\nu\alpha}\right),
\end{equation}
$A$ is the determinant ratio and $\tau$ the corresponding trace defined
in \cite{MagnanoSokolowski}.  The form \eqref{eq:LM} separates Einstein
gravity from a minimally coupled nonlinear massive spin-$2$ field.

For linearly traceless external fields, the determinant and inverse-
determinant contributions to the apparent $\operatorname{tr}\Phi^3$
potential cancel exactly:
\begin{equation}\label{eq:potentialcancel}
V(\Phi)=\frac{m^2}{8}\operatorname{tr}\Phi^2+O(\Phi^4).
\end{equation}
The nonzero $MMM$ amplitude below comes entirely from $K(Q)$.

\subsection{Exact three-point amplitudes}

Take all momenta incoming.  For the $MMM$ vertex use
\begin{equation}\label{eq:mmmkin}
P=(1,0,0,0),\qquad
p=\left(-\frac12,0,0,\frac{i\sqrt3}{2}\right),\qquad
q=\left(-\frac12,0,0,-\frac{i\sqrt3}{2}\right).
\end{equation}
Then $P+p+q=0$ and $P^2=p^2=q^2=-1$.  With normalized TT
polarizations $(M_+,M_+,M_0)$, direct expansion of \eqref{eq:LM} gives
\begin{equation}\label{eq:mmmvalue}
\boxed{\mathcal A_3(M_+,M_+,M_0)=-\frac{\sqrt6}{8}.}
\end{equation}

For $MMh$, take
\begin{equation}\label{eq:mmhkin}
-P=(-1,0,0,0),\qquad
r=(1,-1,-i,0),\qquad
k=(0,1,i,0),
\end{equation}
so $-P+r+k=0$, $r^2=-1$, and $k^2=0$.  For the polarization choice in
Appendix~\ref{app:eframe},
\begin{equation}\label{eq:mmhvalue}
\boxed{\mathcal A_3(M_+,M,h)=-\frac{\sqrt2}{8}.}
\end{equation}
The graviton Ward identity vanishes for an arbitrary symbolic vector
$\xi_\mu$ under
$\varepsilon^h_{\mu\nu}\mapsto k_{(\mu}\xi_{\nu)}$, and exchange of
the two massive legs leaves the coefficient unchanged.

\subsection{Complete massive-pole residue}

Glue \eqref{eq:mmmvalue} and \eqref{eq:mmhvalue} across $P$.  A complete
orthonormal rest-frame basis contains the polarizations
$+,\times,0,xz,yz$.  The exact products are
\begin{center}
\begin{tabular}{c|c|c|c}
\hline
$\lambda$ & $\mathcal A_{MMM}$ & $\mathcal A_{MMh}$ & product\\
\hline
$+$ & $-\sqrt6/8$ & $-\sqrt2/8$ & $\sqrt3/32$\\
$\times$ & $0$ & $0$ & $0$\\
$0$ & $0$ & $\sqrt6/8$ & $0$\\
$xz$ & $0$ & $0$ & $0$\\
$yz$ & $0$ & $i\sqrt2/4$ & $0$\\
\hline
\end{tabular}
\end{center}
Therefore the complete polarization numerator is
\begin{equation}\label{eq:resnum}
N=\sum_{\lambda=1}^{5}
\mathcal A_3(M,M,M_P^\lambda)
\mathcal A_3(M_{-P}^\lambda,M,h)
=\frac{\sqrt3}{32}.
\end{equation}
The TT inverse kernel in the same action convention is
\begin{equation}\label{eq:inversekernel}
K_M(P)=\frac{P^2+M^2}{4}.
\end{equation}
Consequently
\begin{equation}\label{eq:residue}
\boxed{
\Res_{P^2=-M^2}\mathcal A_4(M,M,M,h)=\frac{\sqrt3}{8}\ne0 .}
\end{equation}
An overall rescaling of the action rescales the displayed residue but not
its nonvanishing.  A local quartic contact term is regular at this pole
and cannot cancel it.

\begin{proposition}[Factorization nonvanishing]\label{prop:factorization}
The tree-level four-point amplitude tensor for $MM\to Mh$ is not
identically zero on the relevant complexified on-shell component.
\end{proposition}

\begin{proof}
If the rational amplitude vanished identically, every factorization
residue would vanish.  Equation~\eqref{eq:residue} is the complete
massive-polarization residue and is nonzero.
\end{proof}

\section{The real physical shell}\label{sec:realshell}

\subsection{Open kinematics}

In the center-of-mass frame, with signature $(+,-,-,-)$ restored,
\begin{equation}\label{eq:comkin}
|\mathbf k|=\frac{s-M^2}{2\sqrt s},
\qquad
E_{M,\mathrm{out}}=\frac{s+M^2}{2\sqrt s},
\qquad s\ge4M^2.
\end{equation}
Thus $MM\to Mh$ is physically open on a continuum shell, unlike the
cubic $MMM$ channel.

\begin{lemma}[Complex nonvanishing and a real component]
\label{lem:zariski}
Fix analytic physical polarization sections on an irreducible component
of the complexified four-point mass shell.  The nonsingular real
$MM\to Mh$ region contains a Euclidean-open center-of-mass chart and is
Zariski dense in the complex component containing it.  Hence a rational
tree amplitude that vanishes on an open real subset vanishes identically
on that component and has zero factorization residues.
\end{lemma}

\begin{proof}
Away from exceptional polarization charts, the real physical region is
parametrized by $s>4M^2$, $-1<\cos\theta<1$, and an azimuth.  Its real
points contain an open set of maximal real dimension in the relevant
algebraic component.  If $\mathcal A=P/Q$ is regular there and vanishes
on a real open set, real analyticity makes $P$ vanish there.  The
identity theorem, equivalently Zariski density of that real locus, makes
$P$ vanish on the component.  Then all residues vanish, contrary to
\eqref{eq:residue}.
\end{proof}

The next subsection supplies a stronger direct bridge: an exact nonzero
point already on the real shell.  It makes Lemma~\ref{lem:zariski}
logically redundant for the obstruction theorem, while retaining it as
an independent explanation of why the complex factorization rail is
physically relevant.

\subsection{Interior rational certificate}

Set $M=1$.  With $p_1,p_2$ incoming and $p_3,k$ outgoing, take
\begin{align}\label{eq:rationalpoint}
p_1&=\left(\frac54,0,0,\frac34\right),
&p_2&=\left(\frac54,0,0,-\frac34\right),\nonumber\\
p_3&=\left(\frac{29}{20},-\frac{21}{25},0,-\frac{63}{100}\right),
&k&=\left(\frac{21}{20},\frac{21}{25},0,\frac{63}{100}\right).
\end{align}
Then
\begin{equation}
p_1+p_2=p_3+k,\qquad p_i^2=1,\qquad k^2=0,
\end{equation}
and
\begin{equation}
s=\frac{25}{4}M^2,
\qquad \cos\theta=\frac35.
\end{equation}
This point is interior and noncollinear.  The planar kinematics has a
$y$-reflection selection rule; the verification therefore uses explicit
parity-even real linear polarizations rather than interpreting symmetry
zeros as dynamical cancellations.

\begin{cproposition}[Exact real-shell certificate]\label{cprop:g15}
For the real polarization tensors listed in Appendix~\ref{app:fourpoint},
the complete reduced tree amplitude at \eqref{eq:rationalpoint} is
\begin{equation}\label{eq:AK}
\boxed{\AKred(MM\to Mh)=\frac{7881241032}{5584765625}\ne0.}
\end{equation}
It is real and nonsingular.  The complete sum passes an exact Ward test
for a rational gauge vector, invariance under an exact rational change of
gauge representative, and equality between axial and de Donder internal
gauges.  Hence both $\AKred$ and its normalized Born coefficient $\AK$
are nonzero on a physical-shell neighborhood.
\end{cproposition}

\section{Full fourth-order calculation and the metric obstruction}
\label{sec:full}

The calculation of Proposition~\ref{cprop:g15} is performed directly in
the original fourth-order variables of \eqref{eq:action}, independently
of the Einstein-frame factorization rail.

\subsection{Perturbiner and bordered propagator}

For four plane waves,
\begin{equation}
g_{\mu\nu}=\eta_{\mu\nu}
+\sum_{a=1}^{4}\lambda_a\varepsilon^{(a)}_{\mu\nu}e^{ip_a\cdot x},
\end{equation}
the wave algebra retains only multilinear subsets of the labels and
implements $\partial_\mu\mapsto iP_\mu$ automatically.  The coefficient
of $\lambda_1\lambda_2\lambda_3\lambda_4$ in the action gives the exact
quartic contact.  The same engine produces the quadratic $10\times10$
operator $K(P)$ and cubic current $J_A$ in the symmetric-tensor basis.

No symbolic closed propagator is assumed.  In each $s,t,u$ channel one
solves the exact de Donder-bordered system
\begin{equation}\label{eq:bordered}
\begin{pmatrix}K(P)&C(P)^T\\ C(P)&0\end{pmatrix}
\begin{pmatrix}X\\ \lambda\end{pmatrix}
=\begin{pmatrix}-J^{(cd)}\\0\end{pmatrix},
\qquad
\mathcal A_{\rm exch}=J^{(ab)T}X,
\end{equation}
where
\begin{equation}
C_\mu(X)=P^\alpha X_{\alpha\mu}-\frac12P_\mu X^\alpha{}_\alpha.
\end{equation}
The exact field-equation residual and $CX$ vanish in all three channels.
Replacing $C$ by the axial constraint $n^\mu X_{\mu\nu}=0$ leaves the
total amplitude unchanged.

\subsection{Contact and exchange decomposition}

At the rational point and polarization choice, the four contributions
are
\begin{align}\label{eq:pieces}
\mathcal A_{\rm contact}&=
\frac{8364585636057}{410644531250},\nonumber\\
\mathcal A_s&=-\frac{836906329}{45312500},\nonumber\\
\mathcal A_t&=-\frac{5142800561}{111695312500},\nonumber\\
\mathcal A_u&=-\frac{90855828716}{205322265625}.
\end{align}
Their exact sum is \eqref{eq:AK}.  The contact term alone is not gauge
invariant.  Replacing the external graviton by
$k_{(\mu}\xi_{\nu)}$ makes the \emph{complete} contact-plus-exchange
sum vanish for an arbitrary exact rational test vector; shifting a
physical representative by the same gauge tensor leaves the amplitude
unchanged.  Initial-massive Bose exchange is exact.  The two-point
coefficient against every symmetric-tensor basis direction vanishes for
each external on-shell polarization, excluding equation-of-motion
contamination.

Reversing all external momenta and conjugating the real polarization
sections reproduces the contact and each exchange separately.  The
quadratic kernels are recomputed and satisfy $K(-P)=K(P)$ exactly.  Thus
the reverse process entering the physical adjoint is evaluated, not
inferred.

\subsection{From covariant graphs to the stationary Born operator}

The perturbiner is a covariant Lagrangian calculation, whereas
Proposition~\ref{prop:bornidentity} is stated for the canonical
stationary Born operator.  The following standard tree-level reduction
makes their relation explicit.

\begin{proposition}[Covariant--stationary matching]
\label{prop:covariantborn}
For physical external states at \eqref{eq:rationalpoint}, the complete
contact-plus-exchange coefficient computed above obeys
\begin{equation}
\langle f|\Born_{2,K}(E)|i\rangle
=\eta_{\rm conv}\mathcal N_{fi}\AKred=\AK\ne0,
\end{equation}
where $\Born_{2,K}$ is the standard-real-form version of
\eqref{eq:Bornoperator}.  The same identity holds for the reverse
matrix element.  No absorptive or prescription-dependent term occurs.
\end{proposition}

\begin{proof}
Pass to a gauge-fixed canonical formulation, equivalently to the regular
second-order auxiliary-field form on the split branch.  The stationary
transition operator satisfies the Lippmann--Schwinger expansion
\begin{equation}
\Born_K(E)=V_K+V_K(E-h_0+i0)^{-1}\Born_K(E),
\end{equation}
whose order-$\zeta^2$ coefficient is $v_{2,K}+v_{1,K}G_Ev_{1,K}$.
For each physical mass eigenmode, decomposing the covariant internal
propagator as
\begin{equation}\label{eq:oldfashioned}
\frac{1}{(P^0)^2-\omega_\lambda^2+i0}
=\frac{1}{2\omega_\lambda}
\left(\frac{1}{P^0-\omega_\lambda+i0}
-\frac{1}{P^0+\omega_\lambda-i0}\right)
\end{equation}
turns its two terms into the two old-fashioned time orderings and their
$E-E_n$ denominators.  Summing the complete physical polarization
residues gives the intermediate-state sum in $v_{1,K}G_Ev_{1,K}$,
including the massive residue sign.  The $s,t,u$ graphs exhaust that sum
at four points.

For a derivative or constrained theory, the canonical Legendre transform
can redistribute instantaneous and seagull terms between ``contact'' and
``exchange.''  Their sum is invariant: the canonical quartic matrix
element plus its instantaneous pieces is precisely the four-wave contact
coefficient after the complementary pieces from \eqref{eq:oldfashioned}
are included.  We therefore identify only the complete sum, not the
individual gauge-dependent pieces, with the Born operator.  This is the
usual equivalence of covariant and old-fashioned tree perturbation theory
\cite{WeinbergQFT}; local mass-eigenfield changes preserve its on-shell
matrix element by the equivalence theorem \cite{KOS}.

At the rational point all three bordered internal systems are
nonsingular, and Lemma~\ref{lem:fullcubicshell} excludes an on-shell
cubic intermediate state.  Thus $i0$ in \eqref{eq:oldfashioned} reduces
to the real principal-value resolvent \eqref{eq:reducedresolvent}.  The
external pole residues, box factors, polarization norms and Bose factors
give exactly the real nonzero multiplier in \eqref{eq:LSZmap}.  The
separately evaluated reverse graphs give the ordinary-adjoint reverse
matrix element with the same multiplier.
\end{proof}

\subsection{Quarter-turn and the deformation cocycle}

For the contact, three external massive fields give $(-i)^3=i$.  An
exchange through an internal massless line has endpoint phase
$(-i)^2(-i)=i$.  An exchange through an internal massive line has
endpoint phase $(-i)^3(-i)^2=-i$, while the quarter-turned massive
inverse kernel contributes $-1$, again giving $i$.  Consequently every
piece of \eqref{eq:pieces} obeys Theorem~\ref{thm:phase}, and
\begin{equation}\label{eq:tplus}
\langle Mh|\Born_{2,+}(E)|MM\rangle=i\AK,
\qquad
\langle MM|\Born_{2,+}(E)|Mh\rangle=i\AK.
\end{equation}

Both states have exact free energy
\begin{equation}
E_{MM}=\frac54+\frac54=\frac52
=\frac{29}{20}+\frac{21}{20}=E_{Mh}.
\end{equation}
Let $P_i=|MM\rangle\langle MM|$, $P_f=|Mh\rangle\langle Mh|$, and
$\operatorname{Off}_{if}(X)=P_fXP_i+P_iXP_f$.  On the ordered
orthonormal basis $(|MM\rangle,|Mh\rangle)$,
\begin{equation}\label{eq:Tblock}
\operatorname{Off}_{if}\Born_{2,+}(E)=i\AK\,\sigma_x,
\end{equation}
and hence
\begin{equation}\label{eq:obstruction}
\boxed{
\operatorname{Off}_{if}(P_ES_2P_E)
=-2i\AK\,\sigma_x\ne0.}
\end{equation}
In the reduced external-leg convention, the exact certificate is
\begin{equation}\label{eq:reducedobstruction}
(\eta_{\rm conv}\mathcal N_{fi})^{-1}
\operatorname{Off}_{if}(P_ES_2P_E)
=-\frac{15762482064\,i}{5584765625}\,\sigma_x.
\end{equation}
Equations \eqref{eq:obstruction} and \eqref{eq:reducedobstruction}
follow from the deformation-complex identity \eqref{eq:bornidentity},
rather than from an external-leg phase argument alone.

\subsection{First-order metric ambiguity}

The solution $R_1^{(0)}$ in \eqref{eq:R1spectral} is defined up to a
Hermitian $G$ with $[G,h_0]=0$.  Such a change shifts the second-order
source by $\tfrac12[G,v_1+v_1^{\dg}]$.  The complete regulated
Fock-space statement in Lemma~\ref{lem:fullcubicshell} gives
\begin{equation}\label{eq:ambiguity}
P_E(v_1+v_1^{\dg})P_E=0,
\qquad
P_E[G,v_1+v_1^{\dg}]P_E=0.
\end{equation}
Indeed, $[G,h_0]=0$ makes $G$ block diagonal in the $P_E$ decomposition,
so the second equality is the commutator of $P_EGP_E$ with the first
zero block.  Disconnected spectators and soft configurations have
already been covered by the lemma's regulated Fock-space statement.
The obstruction cannot be moved or canceled by choosing another analytic
first-order metric representative.

\begin{theorem}[Second-order positive-metric obstruction]
\label{thm:positiveobstruction}
The canonical analytic metric deformation based at the Paper-IV
Lorentz-covariant positive real form admits no second-order solution for
the complete Einstein--Weyl tree-level cubic and quartic vertex content.
The obstruction is nonzero on the physical $MM\to Mh$ shell at
\eqref{eq:rationalpoint} and on a nonempty open neighborhood of that
point.
\end{theorem}

\begin{proof}
Propositions~\ref{prop:bornidentity} and \ref{prop:covariantborn} turn
the exact four-point calculation into the nonzero source
\eqref{eq:obstruction} in $\ker\ad_{h_0}$.  No $R_2$ can have a
commutator with $h_0$ equal to that source.  Equation
\eqref{eq:ambiguity} makes the class independent of the homogeneous
first-order metric freedom.  The Ward, external-EOM, reverse-process,
internal-gauge and LSZ arguments show that the displayed matrix element
lies on the reduced physical shell.  Real analyticity extends its
nonvanishing from the interior point to an open neighborhood.
\end{proof}
\section{Conventional Krein grading on physical cohomology}
\label{sec:grading}

The positive obstruction does not by itself decide whether an indefinite
completion can still be formulated.  The narrower question addressed
here is whether the conventional gravitational Krein grading can render
the offending block null, as a continuous one-sided charge does at the
perfect-square scalar boundary.

\subsection{One-particle commutant}

Let $J^{(1)}$ be a nondegenerate Hermitian involution on the physical
one-particle cohomology \eqref{eq:oneparticle}.  We require it to commute
with the proper-orthochronous Poincar\'e representation and to agree with
the free Krein real form of Paper~IV.

The mass Casimir separates the massless and massive orbits, so no
covariant endomorphism mixes $\mathcal H_h$ and $\mathcal H_M$.  On the
massive fiber, the exact spin-$2$ $SO(3)$ commutant is one-dimensional.
There is one precision on the massless fiber: under the
proper-orthochronous group, helicity $+2$ and helicity $-2$ are
inequivalent complex representations, so covariance alone permits
separate scalars.  The complete commutant is therefore
\begin{equation}\label{eq:commutant}
J^{(1)}=\epsilon_+I_{+2}\oplus\epsilon_-I_{-2}
\oplus\epsilon_MI_5,
\qquad \epsilon_\pm,\epsilon_M\in\{\pm1\}.
\end{equation}
Parity, or compatibility with the real gravitational field, exchanges
the two massless helicities and imposes $\epsilon_+=\epsilon_-$.  Even
without adding parity to the covariance group, agreement with the free
signature \eqref{eq:freesignature} fixes
\begin{equation}\label{eq:epsfix}
\epsilon_+=\epsilon_-=+1,
\qquad \epsilon_M=-1.
\end{equation}

\begin{proposition}[One-particle grading classification]
\label{prop:oneparticlegrading}
Within the stated covariant class, the free real form uniquely fixes
\begin{equation}
J^{(1)}=I_2\oplus(-I_5).
\end{equation}
No linear covariant grading can assign different signs within the
massive spin-$2$ multiplet or mix the massive and massless shells.
\end{proposition}

\subsection{Canonical second quantization and cluster multiplicativity}

We do not classify every Poincar\'e-covariant involution on a
multiparticle Hilbert space: such an operator could act on
momentum-dependent multiplicity spaces.  The conventional lift is the
natural second quantization $\Gamma(J^{(1)})$.  Explicitly, we assume
that it is diagonal in asymptotic particle labels, acts on symmetrized
simple tensors by the tensor product of its one-particle actions, and is
multiplicative under separated clusters.  Equivalently, its value on a
sector is a scalar $j(n_h,n_M)\in\{\pm1\}$ depending only on particle
numbers.  These assumptions imply
\begin{equation}\label{eq:monoid}
j(n_h+n'_h,n_M+n'_M)
=j(n_h,n_M)j(n'_h,n'_M),
\qquad j(0,0)=1.
\end{equation}
A homomorphism from the free commutative monoid $\mathbb N^2$ is fixed by
its values on the two generators.  Equations \eqref{eq:epsfix} and
\eqref{eq:monoid} therefore give
\begin{equation}\label{eq:fockgrading}
\boxed{J_F=(-1)^{N_M}.}
\end{equation}
Assigning independent signs to unrelated multiparticle sectors would
violate \eqref{eq:monoid} and hence cluster factorization.

\begin{proposition}[Canonical multiplicative Fock lift]
\label{prop:focklift}
Within the particle-number-diagonal, natural second-quantization class
just specified, $J_F=(-1)^{N_M}$ is the unique lift of
$J^{(1)}=I_2\oplus(-I_5)$.  No claim is made here about arbitrary
momentum-dependent or non-factorizing multiparticle involutions.
\end{proposition}

At the rational physical shell,
\begin{equation}\label{eq:statesigns}
J_F|MM\rangle=+|MM\rangle,
\qquad
J_F|Mh\rangle=-|Mh\rangle.
\end{equation}
The transition crosses the two Krein signatures.

\section{Krein visibility of the obstruction}\label{sec:visibility}

For an operator on the positive-frame Hilbert representatives, define
the Krein adjoint
\begin{equation}\label{eq:sharp}
X^\sharp=J_FX^{\dg}J_F.
\end{equation}
On the minimal ordered basis $(|MM\rangle,|Mh\rangle)$,
$J_{\min}=\operatorname{diag}(1,-1)$.

Consider first the forward block
\begin{equation}\label{eq:Xforward}
X=i\AK\,|Mh\rangle\langle MM|.
\end{equation}
It is $J_F$-odd:
\begin{equation}
J_FXJ_F=-X.
\end{equation}
Its Krein adjoint and quadratic product are
\begin{align}\label{eq:Xsharp}
X^\sharp&=-i\AK\,J_F|MM\rangle\langle Mh|J_F
=i\AK\,|MM\rangle\langle Mh|,\nonumber\\
X^\sharp X&=-\AK^2|MM\rangle\langle MM|.
\end{align}
Consequently
\begin{equation}\label{eq:Xtrace}
\boxed{
\Tr(X^\sharp X)=-\AK^2\ne0.}
\end{equation}
The sign is indefinite, but the operator is not null.

For the off-diagonal positive-metric obstruction block
\begin{equation}
\mathcal O:=\operatorname{Off}_{if}(P_ES_2P_E)=-2i\AK\sigma_x,
\end{equation}
one obtains
\begin{equation}\label{eq:Osharp}
\boxed{
\mathcal O^\sharp=\mathcal O,
\qquad
\mathcal O^\sharp\mathcal O=-4\AK^2I_2,
\qquad
\Tr(\mathcal O^\sharp\mathcal O)=-8\AK^2\ne0.}
\end{equation}

\begin{proposition}[Conventional Krein non-nullity]
\label{prop:nonnull}
The physical $MM\to Mh$ block and its off-diagonal metric obstruction
are non-null for the canonical multiplicative fundamental symmetry
\eqref{eq:fockgrading}.  The obstruction is Krein-self-adjoint,
indefinite, and visible to the finite-block quadratic trace.
\end{proposition}

The algebraic distinction from a one-sided continuous grading is simple.
For $\mathbb Z_2$,
\begin{equation}
1+1=0\pmod2.
\end{equation}
The adjoint of an odd operator is odd and their product is even.  Thus
oddness supplies no charge-null lemma:
\begin{equation}\label{eq:oddnotnull}
\boxed{J_F\text{-odd}\ \not\Longrightarrow\ \text{null}.}
\end{equation}
By contrast, at the scalar $O(1,1)$ boundary a strictly one-sided
nonzero continuous charge remains nonzero after the relevant product,
and an invariant graded trace can kill it \cite{paperV,BT2026}.

\begin{remark}[What is not ruled out]
Equation~\eqref{eq:Osharp} does not say that a Krein-self-adjoint
interacting Hamiltonian cannot exist.  Indeed, it displays
Krein-self-adjointness of the minimal obstruction block.  What fails is
the stronger proposal that the conventional massive-number grading makes
the branch-changing component null or supplies a positive generalized
Born contribution.  Indefinite pseudo-unitary evolution and a positive
probability calculus are different questions.
\end{remark}

\section{No scalar-like charge-null mechanism}\label{sec:charge}

Could a continuous internal charge make the obstruction one-sided?
Poincar\'e covariance first restricts any linear
charge on the irreducible massive fiber to one scalar $q_M$.  The free
real graviton is neutral, $q_h=0$.  The verified nonzero interaction
vertices impose
\begin{equation}\label{eq:charges}
3q_M=0\quad(MMM),
\qquad
2q_M+q_h=2q_M=0\quad(MMh).
\end{equation}

\begin{proposition}[Uniform abelian charge no-go]\label{prop:chargenogo}
There is no nontrivial uniform abelian charge assignment to the physical
massive spin-$2$ irrep, with neutral graviton, that is respected by the
Einstein--Weyl interaction.
\end{proposition}

\begin{proof}
Equations \eqref{eq:charges} hold in any abelian charge group.  Since
$q_M=3q_M-2q_M$, they imply $q_M=0$.  The conclusion does not rely on
torsion-freeness; $\gcd(2,3)=1$ eliminates finite torsion charges as well.
\end{proof}

The $MMM$ vertex has no real cubic decay shell but is a nonzero algebraic
interaction vertex, so it already breaks massive-number parity at the
action level.  More decisively, the exact real-shell process changes
$N_M$ from $2$ to $1$ and proves
\begin{equation}\label{eq:paritybroken}
[(-1)^{N_M},S]\ne0
\end{equation}
on the physical scattering shell.  Therefore the ingredients used to
prove factorization simultaneously exclude an $O(1,1)$-like charge
replacement.

\section{BRST cohomology}\label{sec:brst}

The higher-derivative massive branch is a physical reduced mode and must
not be confused with Faddeev--Popov ghosts.  The external states in the
four-point calculation are TT massive polarizations and TT massless
graviton polarizations.
The exact Ward identity, invariance under changes of gauge representative,
axial/de Donder equality, and external inverse-propagator
zeros show that the amplitude descends to reduced physical BRST
cohomology.

Let $Q$ denote the nilpotent BRST charge and let $|i\rangle,|f\rangle$
be closed physical representatives.  If an operator were BRST-exact,
\begin{equation}
\mathcal O=\{Q,Y\},
\end{equation}
then, under the standard decoupling assumptions,
\begin{equation}\label{eq:brstzero}
\langle f|\mathcal O|i\rangle
=\langle f|QY|i\rangle+\langle f|YQ|i\rangle=0.
\end{equation}
The matrix element \eqref{eq:obstruction} is nonzero.

\begin{proposition}[BRST survival]\label{prop:brst}
The exact obstruction defines a nonzero operator on physical BRST
cohomology and is not BRST-exact:
\begin{equation}
[\mathcal O]\ne0.
\end{equation}
It cannot be placed in the standard BRST-null ideal.
\end{proposition}

This proposition concerns the induced operator on physical cohomology.
We do not construct an extension of $J_F$ to the complete gauge-fixed
Fock space of Faddeev--Popov ghosts and contractible sectors; accordingly
we do not call $J_F$ a full-space BRST-compatible fundamental symmetry.

Appendix~\ref{app:brst} records the finite-complex algebra checked by the
verification suite.
In a splitting into physical cohomology and a contractible doublet, the
physical-to-physical block of $\{Q,Y\}$ vanishes for a completely general
$Y$, while the physical obstruction block survives unchanged.

\section{Master theorem}\label{sec:master}

\begin{theorem}[Interacting completion obstruction]\label{thm:master}
Let scalar-free Einstein--Weyl gravity \eqref{eq:action} be in its regular
split phase about flat space.  Use the covariant free completion
classes of Paper~IV and the analytic metric-deformation class of
Paper~V.  Then:
\begin{enumerate}
\item the Lorentz-covariant positive pseudo-Hermitian completion is
protected at first perturbative order, $\ob_+^{(1)}=0$;
\item at second order the physical process $MM\to Mh$ generates the
nonzero shell obstruction \eqref{eq:obstruction};
\item hence the canonical Paper-IV positive base metric has no analytic
second-order deformation for the complete tree-level cubic and quartic
Einstein--Weyl vertex content;
\item on physical BRST cohomology, the free one-particle real form has
the canonical particle-number-diagonal, natural and
cluster-multiplicative lift $J_F=(-1)^{N_M}$;
\item for this grading the obstruction is non-null,
$\Tr(\mathcal O^\sharp\mathcal O)=-8\AK^2\ne0$, and survives physical
BRST cohomology; and
\item the nonzero $MMM$ and $MMh$ vertices prohibit a nontrivial uniform
abelian charge grading analogous to the scalar $O(1,1)$ mechanism.
\end{enumerate}
Thus, through second tree order, neither the canonical analytic metric
deformation nor the conventional massive-number null mechanism supplies
a positive-probability completion in the stated classes.
\end{theorem}

\begin{proof}
Item 1 is Theorem~\ref{thm:cubicprotection}.  Items 2 and 3 are
Theorem~\ref{thm:positiveobstruction}, independently supported by the
factorization Proposition~\ref{prop:factorization} and the
Born--deformation identity \eqref{eq:bornidentity}.  Item 4 is
Propositions~\ref{prop:oneparticlegrading} and \ref{prop:focklift}.
Item 5 is
Propositions~\ref{prop:nonnull} and \ref{prop:brst}.  Item 6 is
Proposition~\ref{prop:chargenogo}.
\end{proof}

\section{Scope and nonclaims}\label{sec:scope}

The hypotheses in Theorem~\ref{thm:master} are part of the result.  The
paper rules out:
\begin{itemize}
\item analytic deformation through second order of the canonical
Paper-IV Lorentz-covariant positive base metric by the complete
tree-level cubic and quartic vertex content;
\item a conserved naive massive-number ghost parity;
\item conventional covariant linear one-particle fundamental symmetries
that agree with the free real form but differ on selected massive
helicities;
\item independent particle-number-sector signs within the natural,
cluster-multiplicative scalar Fock-lift class;
\item a uniform abelian continuous- or torsion-charge null mechanism
compatible with the verified vertices; and
\item standard BRST-exact nullity of the physical obstruction.
\end{itemize}

It does \emph{not} prove:
\begin{itemize}
\item that no nonperturbative or differently pointed positive completion
exists;
\item that the full interacting Hamiltonian has complex energies;
\item that every nonlocal or non-factorizing grading fails;
\item that no momentum-dependent Poincar\'e-covariant involution exists
on multiparticle multiplicity spaces;
\item that the cohomological grading $J_F$ necessarily extends to a
fundamental symmetry on the complete gauge-fixed BRST complex;
\item that no singular $M\to0$ or conformal-boundary completion can be
defined;
\item that a Maldacena-type boundary truncation of the solution algebra is
inconsistent \cite{Maldacena}; or
\item that the result extends unchanged to quadratic gravities with a
scalaron, matter, cosmological constant, or a different background.
\end{itemize}

Two prominent prescriptions lie outside the assumptions for a
particularly concrete reason.  In the unstable-resonance treatment of
Donoghue and Menezes \cite{DonoghueMenezes}, the massive ghostlike pole
does not belong to the asymptotic spectrum and only stable states enter
the unitarity sum.  In the fakeon construction
\cite{AnselmiFakeon}, the massive mode is purely virtual and amplitudes
are defined by a nonanalytic continuation.  The present theorem instead
puts physical $M$ states on both sides of the scattering process and
studies an analytic metric deformation.  It therefore neither refutes
nor subsumes either prescription.

Any surviving null-sector construction must therefore leave the
conventional class.  Possibilities include a grading singular at the
conformal boundary, a nonlinear mixing of asymptotic particle sectors, a
nonlocal observable algebra, a failure of ordinary cluster
factorization, or boundary conditions that remove part of the
fourth-order solution space.  These are qualitatively different from a
straightforward gravitational lift of the Bateman--Turok mechanism.

There is a second, logically independent reason that such a lift cannot
be called straightforward.  Paper~V \cite{paperV} proves, with dependency tag
\texttt{REDUCED-MODE}, that the complete
final-state collinear real normalization response does not cancel the virtual
external-mass logarithm under any continuous axis-compatible parent/daughter
gluing, including the physical pair threshold.  This does not transfer the
scalar coefficient or obstruction to gravity.  It identifies an additional
construction that a gravitational completion would need independently: a
distributional normalization, degenerate or dressed asymptotic states, or a
resummation prescription defined on the gravitational BRST carrier.
The scalar certificate
\texttt{REVERSE\_PHYSICS\_BT\_ASYMPTOTIC\_}\allowbreak
\texttt{GENERATOR\_PREFLIGHT\_V1}
gives an exact finite witness for the dressed alternative.  After division
by the scalar Born rate, projector idempotence forces dimensionless hard and
three-pair blocks $-1/16$ and $+1/16$; multiplying back gives the physical
$-3/512$ normalization against the scalar real response.  The same certificate
shows that the ordinary single-denominator cubic Fock generator vanishes in
the massless collinear limit, so a Jordan/distributional asymptotic term is
required.  The scalar certificate
\texttt{REVERSE\_PHYSICS\_BT\_RT\_}
\allowbreak\texttt{JORDAN\_KERNEL\_V1} derives the resonant
two-annihilator part of that term from the published composite map.  After a
declared repair of an internal Appendix~C oscillator-label inconsistency, its
apparent $t,t^2$ growth cancels exactly on the $(\Omega,\Upsilon)$ carrier and
a formal finite-mode anti-Krein generator exists.  The resulting Krein Gram
has cubic splitting-endpoint poles, so its distributional continuum domain,
oscillatory sectors, and the target $1/48$ remain open.  Gravity would still
require an independent BRST-compatible construction rather than importing
either scalar coefficients or this reduced-mode endpoint obstruction.
The scalar endpoint extension retains three
reflection-even local constants; distinct exact plus and cutoff prescriptions
disagree even on the inclusive constant test.  Hence neither scalar symmetry
nor the interior kernel selects the target coefficient without the omitted
oscillatory and squeezed sectors, and no such ambiguity may be imported as a
gravitational normalization.
On the published fixed-vacuum oscillator grading, the scalar oscillatory and
squeezed sectors lie in the strictly negative BT charge radical and are
trace-null.  The covariant vacuum-orbit calculation shows that this
exclusion does not transfer to the broken-vacuum carrier.  Conversely, the
free cross-CCR and projection idempotence
force the scalar coisometric support to be the identity to every formal order.
The remaining scalar canonical identities are insufficient selectors.  The
exact certificate
\texttt{REVERSE\_PHYSICS\_BT\_CANONICAL\_}
\allowbreak\texttt{ENDPOINT\_AMBIGUITY\_V1}
constructs a three-parameter canonical/projector countermodel in the endpoint
carrier: hard and endpoint trace shifts cancel for arbitrary parameters, and
$1/48$ is compatible but not derived by those identities.  This is not a
claim that all parameters
occur in the actual BT transport; it isolates the deferred continuum map or
an independent physical endpoint condition as the missing scalar datum.  It
is not a range defect or phase prescription that gravity could copy.
The scalar certificate
\texttt{REVERSE\_PHYSICS\_BT\_FULL\_OFF\_}
\allowbreak\texttt{RESONANT\_PROJECTOR\_V1} retains the exact
parent--daughter energy deficit and restores the full radial daughter measure.
The repaired carrier has no surviving $t,t^2$ powers and restricts exactly to
the resonant kernel, but its universal soft cross-Gram is $-1/(2r^3)$.
Multiplication by $r^2dr$ leaves $-(1/2)dr/r$, whose finite part shifts by
$(1/2)\log c$ under a common cutoff rescaling.  Thus full phase space reduces
the collinear three-jet ambiguity to one soft normalization on the declared
chart without selecting it.  A scalar soft-collinear asymptotic Hamiltonian
or hard matching operator with cancelling regulator response is required.
This is only a \texttt{LOCAL-ALGEBRAIC}, \texttt{REDUCED-MODE}
classification; it is not a graviton/BRST transfer, a complete probability,
or a \texttt{LORENTZIAN-CAUSAL} result.
The scalar certificate
\texttt{REVERSE\_PHYSICS\_BT\_SOFT\_CHARGE\_}
\allowbreak\texttt{RESOLVED\_FLOW\_V1} fixes the numerical normalization but
exposes a charge obstruction.  Parent-index raising, Bose reduction,
angular measure, and the outgoing $2!/3!$ projector ratio yield exactly
$+1/48$ per unordered pair and $+1/16$ in total; idempotence forces the hard
$-1/16$.  Yet the two logarithmic scalar rows have generator charges
$(+1,-1)$ and $(-1,+1)$.  Projecting the published oscillator map first to
its one-sided nonpositive image therefore gives zero, while obtaining the
candidate neutral $1/48$ requires a background zero-mode completion carrying
both compensating signs.  That operator and the full order-$\lambda$
pushforward are not public.  This trichotomy is
\texttt{LOCAL-ALGEBRAIC}, \texttt{REDUCED-MODE}; none of its scalar charge
grading, coefficient, or formal finite-cutoff anti-Krein flow transfers to
the metric BV complex.
The exact scalar certificate
\texttt{REVERSE\_PHYSICS\_BT\_ZERO\_MODE\_}
\allowbreak\texttt{EQ19\_TRILEMMA\_V1} constructs the global shift-orbit
operator $Z=e^{\lambda\phi_0}$.  Equation~(16) factorizes as
$\Omega=Z\widehat\Omega$, $\Upsilon=Z^{-1}\widehat\Upsilon$, and the unique
covariant dressing neutralizes both soft logarithmic rows while leaving their
conditional coefficient at $1/48$.  But it also changes the Appendix-C
squeeze to neutral $Z^2(b_\Upsilon^\dagger)^2$.  Fixing $Z=1$ recovers the
apparent charge $-2$, yet the ideal $(Z-1)$ is not invariant because
$\delta(Z-1)=Z\equiv1$ modulo that ideal.  The deferred full zero-mode module,
trace, and nonlinear pushforward are therefore indispensable.  This is a
scalar \texttt{LOCAL-ALGEBRAIC}, \texttt{REDUCED-MODE} obstruction only; it
neither establishes Eq.~(19) nor transfers a coefficient, state, positivity
claim, or causal construction to gravity.
The further scalar certificate
\texttt{REVERSE\_PHYSICS\_BT\_SQUEEZED\_}
\allowbreak\texttt{VACUUM\_IMPLEMENTABILITY\_V1} tests the displayed
Appendix-C vacuum relation in the positive topology of its ordinary carrier.
Its normalized creation kernel is $1/(8|\mathbf p|^2)$; the first
two-particle norm density is
$\rho^2(\epsilon^{-1}-\Lambda^{-1})/(64\pi^2)$, and a six-mode lower bound
grows linearly with box size for every uniformly equivalent fundamental
symmetry.  The pair block independently fails the Hilbert--Schmidt condition.
Hence Krein nullity is not ordinary Fock normalizability.  This exact scalar
\texttt{LOCAL-ALGEBRAIC}, \texttt{REDUCED-MODE} obstruction leaves extended
or inequivalent non-Fock representations open and supplies no graviton state,
BRST theorem, coefficient, or \texttt{LORENTZIAN-CAUSAL} result.
The full scalar exponential can be made modewise convergent by the explicit
weight
$\rho_\mu(p)=4\gamma\mu^2p^2/(p^2+\mu^2)$, $0<\gamma<1$, but its inverse
is unbounded and the thermodynamic vacuum is inequivalent to the ordinary
Fock sector.  The raw scalar shear also fails the positive-Hilbert Bogoliubov
relation, so generic extended positive-boson implementation theorems do not
supply its cross-Krein map or cyclic Born trace.  A direct scalar construction
does implement the covariant squeeze on paired cores and preserves the
finite-rank trace.  It also proves the limiting obstruction: a finite
normalized cyclic orbit trace positive on the ghost-even cone assigns zero to
localized rank-one projectors, whereas the finite-rank trace has infinite
identity weight, and the transported vacuum projector's trace norm grows as
$\exp(V\ell)$ with
\[
 \ell=\frac{\mu^3}{12\pi}\left[
 (1+\gamma)^{3/2}+(1-\gamma)^{3/2}-2\right]>0.
\]
Thus no normal thermodynamic Born trace follows.  This remains solely a scalar
carrier classification.  Its dependency tags are \texttt{LOCAL-ALGEBRAIC}
and \texttt{REDUCED-MODE}; it does not construct a graviton, BRST, or causal
state space.
A scalar semifinite construction retains localized orbit projectors without
assigning finite weight to the identity.  Its canonical trace satisfies
$\operatorname{Tr}_{\!\infty}(E_n)=1$ and
$\operatorname{Tr}_{\!\infty}(1)=+\infty$.  For a finite incoming Krein
projection and a finite exhaustive output partition, the relative weights
\[
 p_i=\frac{\operatorname{Tr}_{\rm fin}(A_i^\dagger A_i)}
 {\operatorname{Tr}_{\rm fin}(P_{\rm in})}
\]
are nonnegative and sum to one whenever each process $A_i$ satisfies the BT
weak ghost decomposition.  The normalized corner functional is not a trace:
$P=E_{00}$, $X=E_{01}$, $Y=E_{10}$ give
$\omega_P(XY)=1$ and $\omega_P(YX)=0$.  Hence this result preserves the
finite-normalized-trace no-go while constructing a finite-detector conditional
calculus.  The exact scalar certificate
\texttt{REVERSE\_PHYSICS\_BT\_SEMIFINITE\_}
\allowbreak\texttt{RELATIVE\_BORN\_WEIGHT\_V1} carries dependency tags
\texttt{LOCAL-ALGEBRAIC} and \texttt{REDUCED-MODE}.  It does not control the
unbounded squeeze on the full trace ideal, construct a thermodynamic state,
reproduce Eq.~(19), transfer to the metric BV complex, or establish any
\texttt{LORENTZIAN-CAUSAL} statement.
The complete public scalar quadratic map closes over three leading inverse
images for every target $b_\Upsilon(s,p)$.  The off-resonant extraction phase
of the opposite-sign $a_2(-s,-p)$ image exactly cancels its Appendix-C phase.
The signed sum is time independent and has only
\[
 \delta b_\Omega=
 \frac{s_1e_1+s_2e_2}{2e_1e_2}
 b_\Omega^{(s_1)}b_\Omega^{(s_2)},
 \qquad
 \delta b_\Upsilon=-\frac{s_2}{2e_1}
 b_\Omega^{(s_1)}b_\Upsilon^{(s_2)}
 -\frac{s_1}{2e_2}
 b_\Upsilon^{(s_1)}b_\Omega^{(s_2)}.
\]
Its parent cross-Gram vanishes, so the resonant-only $dr/r$ carrier and its
conditional $1/48$ normalization do not occur in the completed quadratic
map.  The mixed-sign rows satisfy the exact order-$\lambda$ cross-CCR Ward
identity; after the covariant $Z$ dressing, a neutral finite-mode projector has
$R_tP_0R_t^\dagger=P_0+\lambda[K,P_0]+O(\lambda^2)$ with $Q_1=0$ in the
Eq.~(19) decomposition.  The certificate
\texttt{REVERSE\_PHYSICS\_BT\_FULL\_SIGNED\_}
\allowbreak\texttt{QUADRATIC\_CLOSURE\_V1} has dependency tags
\texttt{LOCAL-ALGEBRAIC} and \texttt{REDUCED-MODE}.  This Eq.~(19)
pushforward is not the scalar S-matrix transition block and establishes no
metric-BV or \texttt{LORENTZIAN-CAUSAL} statement; its
squeezed-vacuum/dynamical-zero-mode projector trace, continuum domain, and
higher orders remain open.
On the finite paired core, cross-Krein isometry gives
$(SKS^{-1})^\dagger(SKS^{-1})=S(K^\dagger K)S^{-1}$ while the detector is
transported by the same similarity.  The cyclic finite-rank trace therefore
keeps the completed scalar quadratic coefficient at zero.  A bare one-pair
occupation of the squeezed vacuum compares a fixed $b$-Fock detector and is
not a further Eq.~(19) term.  The scalar certificate
\texttt{REVERSE\_PHYSICS\_BT\_SQUEEZED\_}
\allowbreak\texttt{DETECTOR\_SIMILARITY\_V1} records this
\texttt{LOCAL-ALGEBRAIC}, \texttt{REDUCED-MODE} distinction without promoting
it through the non-trace-class thermodynamic limit or into the metric BV
complex.
For the physical scalar process, the independently computed five-point
response remains $+1/512$ per pair and $+3/512$ over three pairs, equivalently
$1/48$ and $1/16$ after Born normalization.  The axis-compatible physical
virtual daughter-ratio response is zero, leaving $+3/512$ uncancelled.  The
public $R_t$ response is also zero but is not an S-matrix summand.  The
object-typed certificate
\texttt{REVERSE\_PHYSICS\_}
\allowbreak\texttt{BT\_INCLUSIVE\_NLO\_}
\allowbreak\texttt{OBJECT\_LEDGER\_V1} therefore requires a separate
hard/dressed response $-3/512$ and leaves the complete scalar NLO probability
unestablished.  None of this scalar ledger transfers to the metric BV complex.
A scalar conditional completion theorem further shows that any regulated
physical shell S-matrix with
$S_{\rm phys}=1+xA+x^2B+\cdots$ and $S_{\rm phys}^\dagger S_{\rm phys}=1$
must obey $2\operatorname{Re}B_{hh}=-\|Ah\|^2=-1/16$.  It therefore forces
the absolute hard response $-3/512$ from the measured real column, independently
of phases.  The exact finite witness in
\texttt{REVERSE\_PHYSICS\_BT\_PHYSICAL\_SHELL\_}
\allowbreak\texttt{PSEUDOUNITARY\_COMPLETION\_V1} does not construct the
continuum dressed scalar S-matrix, prove beyond-tree positivity, or supply a
metric-BV carrier.
The scalar logarithmic-shell limit sharpens this boundary.  Disjoint endpoint
shells have $\|A_nh-A_mh\|^2=1/8$, and their exact isometries have hard-column
distance squared $2\sin^2(x/4)$, excluding an ordinary strong Møller limit.
The weak limit loses the real channel and is not pseudo-unitary.  The exact
fixed endpoint-fibre pullback certified in
\texttt{REVERSE\_PHYSICS\_BT\_LOG\_SHELL\_}
\allowbreak\texttt{MOLLER\_LIMIT\_V1}
preserves the leading-log scalar probability cylinder but supplies neither a
local detector affiliation nor a metric-BV construction.
The scalar endpoint fibre is now affiliated with an asymptotic
mass-resolution detector algebra: for a monotone profile $q_R(y)=q(y-R)$,
the positive difference $d_{R,a}=q_{R+a}-q_R$ has exact trace $a$, and its
normalized square root is translation covariant.  Sharp and $C^1$ cubic
profiles therefore give the same $+1/16$ real and $-1/16$ hard response.
Certificate
\texttt{REVERSE\_PHYSICS\_BT\_DETECTOR\_}
\allowbreak\texttt{RESOLUTION\_DILATION\_V1}
establishes this physical scalar leading-log resolution response only on its
declared final-pair cylinder.  Abel integration of the formal scalar soft
Hamiltonian gives coefficient $-id/(\epsilon+id)$ and logistic norm profile
$q_R(y)=(1+e^{2(y-R)})^{-1}$ with $R=\log(\alpha/\epsilon)$.  Thus inverse
Abel time implements the same resolution translation.  The coherent columns
have fixed-ratio squared distance
$\log(c)(c-1)/(c+1)>0$; the resulting orthogonal-increment Naimark carrier is
certified in
\texttt{REVERSE\_PHYSICS\_}
\allowbreak\texttt{BT\_ABEL\_NAIMARK\_}
\allowbreak\texttt{ASYMPTOTIC\_}
\allowbreak\texttt{DILATION\_V1}.
This identifies a scalar detector automorphism, not the formal $R_t$ lowering
operator with the physical S-matrix splitting operator.  Its auxiliary
resolution coordinate, scalar detector trace, cancellation, and coefficient
do not transfer to the metric BV complex or establish a
\texttt{LORENTZIAN-CAUSAL} result.
The scalar amplitude-level comparison has also been completed on its declared
external mass-jet cylinder.  The physical splitting map has full-rank raised
Gram $\rho I_2$ above unequal-mass threshold, whereas the completed public
quadratic $R_t$ map has rank-one nilpotent raised Gram
$\left(\begin{smallmatrix}0&2\\0&0\end{smallmatrix}\right)$.  The exact
rank/Jordan obstruction in
\texttt{REVERSE\_PHYSICS\_BT\_PHYSICAL\_}
\allowbreak\texttt{COLLINEAR\_OPERATOR\_}
\allowbreak\texttt{FACTORIZATION\_V1}
prevents a scalar metric-compatible identification of those two operators;
it neither supplies a metric-BV splitting map nor transfers the scalar
$1/48$, phase, carrier, or positivity statement to gravity.
The scalar map also has a finite-regulator direct-integral realization, but
its exact Gram
$I(r)=5/24+r[(1/4)\log r+1/12]+O(r^2\log r)$ is not differentiable at the
massless axis.  Hence no strongly $C^1$ scalar M\o ller column on a fixed
bounded-pairing carrier can implement the delta-prime derivative.  The
relative $1/48$ scale cocycle is linearized only on the rigged affine sector
$\operatorname{span}\{1,s\}\subset\mathcal S'(\mathbb R_s)$, whose
generators are not $L^2$ vectors.  Certificate
\texttt{REVERSE\_PHYSICS\_BT\_RIGGED\_}
\allowbreak\texttt{RESOLUTION\_JORDAN\_}
\allowbreak\texttt{MOLLER\_V1}
does not supply a scalar endpoint state or generalized-Born functional, and
none of its rigging, cocycle, or differentiability statement transfers to the
metric BV complex.
On the scalar resolution-local CCR net, the rank-two Gram further has a
positive coherent realization.  A finite interval of length $a$ has coherent
norm $a/16$, hard no-emission probability $e^{-a/16}$, and normalized Poisson
count law $e^{-a/16}(a/16)^n/n!$.  The compatible local Weyl implementers
define a locally normal algebraic state, while the norm grows as $L/16$ and
excludes a global scalar Fock implementer.  The certificate
\texttt{REVERSE\_PHYSICS\_BT\_RESOLUTION\_}
\allowbreak\texttt{LOCAL\_COHERENT\_BORN\_}
\allowbreak\texttt{PROCESS\_V1}
constructs this scalar leading-log process only under a coherent independent-
increment completion.  It neither derives scalar multi-emission dynamics nor
supplies a spacetime-local scalar S-matrix, and no part of its auxiliary CCR
net, Poisson law, or rank-two scalar GNS factor transfers to metric BV data.
The subsequent complete scalar six-point tree calculation rules out that
independent-increment completion in its nested strongly ordered sector:
$P_2(a)=5a^2/512$ rather than $a^2/512$, with scalar second factorial
cumulant $a^2/64$.  This factor-five result is certified by
\texttt{REVERSE\_PHYSICS\_BT\_SIX\_POINT\_}
\allowbreak\texttt{STRONGLY\_ORDERED\_TREE\_V1}; it requires a non-Gaussian
scalar resolution state but does not select one from only two moments.  The
calculation uses scalar external mass jets and scalar K\"all\'en thresholds.
Neither its coefficient nor its non-Poisson interpretation transfers to the
metric BV carrier.
The complete scalar seven-point tree now supplies the third ordered moment:
$P_3(a)=9a^3/8192$ and $\kappa_3(a)=7a^3/2048$.  It excludes the gamma--Cox
completion, which predicts $15a^3/8192$, while satisfying the scalar
Stieltjes moment bound.  The certificate
\texttt{REVERSE\_PHYSICS\_BT\_}
\allowbreak\texttt{SEVEN\_POINT\_}
\allowbreak\texttt{COX\_SELECTION\_V1}
constructs a positive locally normal two-atom Cox state matching the first
three scalar moments.  This state is an auxiliary resolution-local controlled
Weyl mixture, not a scalar spacetime M\o ller operator, a complete scalar
probability, or an all-order dynamical law.  Its coefficient, intensity
carrier, positivity construction, and threshold calculation do not transfer
to the metric BV complex.  The complete scalar eight-point preflight does not
yet add a fourth moment.  All $34300$ trees give a bare projected hierarchy
of ordered valuation $(0,0,-1)$, and two fixtures with identical soft data
but one hard adjacent invariant changed from $33$ to $34$ have leading
residues differing by exactly $257/1568$.  Certificate
\texttt{REVERSE\_PHYSICS\_BT\_EIGHT\_POINT\_}
\allowbreak\texttt{HARD\_PROFILE\_OBSTRUCTION\_V1}
therefore leaves the threshold-integrated scalar fourth moment and the
two-atom Cox test open until the outer K\"all\'en reduction or a hard-profile
quotient is computed.  The exact outer reduction retains the two full
profiles and finds their leading difference to be the pure $J_2$ direction
$771/(1568e_3u)$.  Since the fixed-invariant $r\log r$ coefficient of $J_2$
is one, the threshold coefficient remains nonzero, $771/(1568e_3)$; at the
declared $r=5e_3$ the outer logarithms differ by $3855/1568$.  Certificate
\texttt{REVERSE\_PHYSICS\_BT\_EIGHT\_POINT\_}
\allowbreak\texttt{OUTER\_THRESHOLD\_PROFILE\_V1}
is followed by the symbolic $\tau_3$ profile.  Its exact two-fixture
difference is
$3(27213\tau_3^2-13462\tau_3+29485)/(156800\tau_3^2\tau_4)$.
The next $J_1,J_2,J_3$ reduction at $a_3=2$ leaves
$334239/313600\ne0$.  Certificate
\texttt{REVERSE\_PHYSICS\_BT\_EIGHT\_}\allowbreak
\texttt{POINT\_MIDDLE\_THRESHOLD\_}\allowbreak
\texttt{PROFILE\_V1}
is followed by the symbolic $\tau_2$ profile.  The two complete hard fixtures
differ by
$3(2090\tau_2^2+1146\tau_2+981)/(12800\tau_2^2)$, and the third
$J_1,J_2,J_3$ reduction at $a_2=3$ leaves $23229/6400\ne0$.
Certificate
\texttt{REVERSE\_PHYSICS\_BT\_EIGHT\_}\allowbreak
\texttt{POINT\_TAU2\_THRESHOLD\_}\allowbreak
\texttt{PROFILE\_V1}
leaves only the independent-mass inner threshold open.  The displayed value
tests profile collapse; it is not a normalized fourth moment.  None of the
scalar profile certificates is imported into the metric BV complex, and none
supplies a gravitational fourth jump.
The independent-mass test fixes only the redundant scale $a_0=1$ and retains
$r=a_1/a_0$, $u=\tau_1/a_0$.  The complete profile difference is
$3\{9r^2-18r(u+1)+34u^2-18u+9\}/(128u^2)$.  Exact invariant-cutoff residues
give the two $r\log r$ coefficients $-6699/128$ and $-7149/128$, whose
difference is $225/64\ne0$.  Certificate
\texttt{REVERSE\_PHYSICS\_BT\_EIGHT\_}\allowbreak
\texttt{POINT\_INNER\_THRESHOLD\_}\allowbreak
\texttt{OBSTRUCTION\_V1}
therefore obstructs the hard-independent scalar fourth jump on the declared
fixtures after all four threshold functionals.  A profile/channel-valued
successor remains open; no normalized probability or gravitational result is
inferred.
This profile direction cannot be identified with the forced cross-Krein
complement alone.  The two fixtures share $\rho=819/4000$, the same missing
Gram and $c_1=14819/8000$, yet have unequal fourth responses.  Certificate
\texttt{REVERSE\_PHYSICS\_BT\_COMPLEMENT\_}\allowbreak
\texttt{HARD\_PROFILE\_SEPARATION\_V1}
requires at least one additional hard-profile base coordinate beyond $\rho$.
It fixes an algebraic type, not a BT dynamical bundle or Eq.~(19).
On the resulting two-point base, the fourth effect is
$K_4=\operatorname{diag}(-6699/128,-7149/128)$.  Its external eight-leg sign
and every inherited normalization sign are positive, so every faithful
positive profile trace is negative.  Certificate
\par\noindent
\texttt{REVERSE\_PHYSICS\_BT\_EIGHT\_}\allowbreak
\texttt{POINT\_PROFILE\_POSITIVITY\_}\allowbreak
\texttt{OBSTRUCTION\_V1}
\par\noindent
therefore obstructs a diagonal completely positive fourth jump.  It also
constructs the minimal two-negative-direction Krein-skew lift over the hard
base, but not a positive probability, charge-compatible BT operator, or
Eq.~(19).
The invariant charge split further shows that this lift is not the
negative-charge $Q$ remainder.  Transporting the actual target charge action
through the certified missing leg $C$ gives the null charge basis of
$G_{\rm miss}$, in which
the vector used by the lift has two separately null charge components; its
norm is entirely their cross pairing.  Thus
$B_+^\sharp B_+=B_-^\sharp B_-=0$ while
$B_+^\sharp B_-+B_-^\sharp B_+=K_4$.  Certificate
\par\noindent
\texttt{REVERSE\_PHYSICS\_BT\_EIGHT\_POINT\_}\allowbreak
\texttt{KREIN\_CHARGE\_LOCALIZATION\_V1}
\par\noindent
localizes the effect in the total-charge-zero sector and leaves the neutral
higher-composite or dynamical trace unconstructed.
The smallest neutral symmetric-bosonic carrier can now be classified.
Over the two hard-profile modes its degree-two charge-zero sector has inertia
$(3,1)$ and cannot carry the rank-two negative profile block, whereas degree
four has inertia $(6,3)$.  Two normalized occupation-antisymmetric degree-four
vectors are total-charge zero, mutually orthogonal and have Gram $-2I_2$;
the certified amplitudes therefore reconstruct $K_4$ in a charge-compatible
Krein-skew generator.  Certificate
\par\noindent
\texttt{REVERSE\_PHYSICS\_BT\_NEUTRAL\_}\allowbreak
\texttt{BOSONIC\_COMPOSITE\_LIFT\_V1}
\par\noindent
also shows why this is not yet Eq.~(19): the negative neutral subspace is odd
under the induced ghost parity, while the positive profile source is even.
Every ghost-even forward block then has pullback
$D^T(G_4\kappa)D\succeq0$, so it cannot reproduce negative $K_4$.
Thus the minimal block is not by itself the ghost-even neutral $P$ term; a ghost-odd
source partner or larger dynamical projector/trace remains necessary.
The minimal two-dimensional odd partner does produce an exact finite
projector.  With partner metric $-I_2$, composite metric $-2I_2$, slope
$T=\operatorname{diag}(\sqrt{6699}/16,\sqrt{7149}/16)$ and
$L=(I_2,T)^T$, the graph projector
$P=L(L^T\eta L)^{-1}L^T\eta$ is rank two, idempotent, Krein self-adjoint,
neutral and ghost even, with finite algebraic trace
$\operatorname{tr}(P^\sharp P)=2$.  Certificate
\par\noindent
\texttt{REVERSE\_PHYSICS\_BT\_NEUTRAL\_}\allowbreak
\texttt{GRAPH\_PROJECTOR\_EXTENSION\_V1}
\par\noindent
also proves that this minimal graph cannot be affiliated to the original
positive even profile source: parity forces the direct map to vanish and
the graph metric is negative definite.  A larger signature-balanced
$R_t$ projector or dynamical trace remains necessary.
The required odd source type is nevertheless already present in the full
public carrier.  In the two-profile neutral degree-four $J$-Fock sector,
which has inertia $(6,3)$, the two occupation-antisymmetric vectors normalized
by $1/2$ have Gram $-I_2$, charge zero and odd ghost parity.  Since the
certified complement map obeys $CS=-\rho I_2$ on the charge basis,
$\operatorname{Sym}^4(C)=\rho^4I_9$ and maps the complement vectors to
$\sqrt2$ times these public vectors.  Its inverse $I_2/\sqrt2$ is therefore
an exact ghost-even isometry from the public odd sector to the selected
composite.  Certificate
\par\noindent
\texttt{REVERSE\_PHYSICS\_BT\_PUBLIC\_FOCK\_}
\allowbreak\texttt{ODD\_SOURCE\_AFFILIATION\_V1}
\par\noindent
replaces the abstract odd summand by a concrete public Fock subspace.  It
does not affiliate the original positive scalar source or establish a
physical fourth probability.  The graph slope belongs to the eight-point
physical-response ledger, not automatically to the distinct
$R_tP_\chi R_t^\dagger$ projector ledger; their identification requires an
explicit intertwiner.
On the covariant zero-mode Laurent algebra the charge-support part of
Eq.~(19) nevertheless holds to all formal orders.  The exact Eq.~(16)
pullback sends $\Omega$ to
$\lambda^{-1}Ze^{\lambda\varphi}$ and $\Upsilon$ to
$Z^{-1}e^{-\lambda\varphi}
[\Box\varphi+\lambda(\partial\varphi)^2]$, so it intertwines the source
orbit derivation and target boost charge coefficient by coefficient.  Free
time translation preserves the identity because both free Hamiltonians
are charge neutral.  Perturbative
two-sidedness makes the inverse pushforward equivariant.  Therefore every
finite shift-invariant nonzero-mode projector remains idempotent,
self-adjoint and wholly charge zero, with $Q_{\rm negative}=0$.  Certificate
\par\noindent
\texttt{REVERSE\_PHYSICS\_BT\_COVARIANT\_FORMAL\_}
\allowbreak\texttt{EQ19\_CHARGE\_SUPPORT\_V1}
\par\noindent
does not promote the full Eq.~(19): ghost evenness, time independence,
asymptotic limits and the continuum trace remain open, and the charge theorem
does not descend through $Z=1$ because
$\delta_\phi(Z-1)=Z\equiv1\pmod{(Z-1)}$.
Ghost parity supplies a separate obstruction already at first nonlinear
order.  The complete public signed kernel has a charge-$+1$ Krein-skew
species generator $K_+$, and covariance fixes the neutral generator to
$Z^{-1}K_+$.  For the standard two-species one-particle projector,
$P_1=Z^{-1}[K_+,P_0]$ is nonzero, rank four and stationary, whereas
$\kappa P_1\kappa=Z\kappa[K_+,P_0]\kappa$.  Independence of the two Laurent
powers proves that the neutral correction is not ghost even.  Its canonical
odd remainder is orthogonal to the even part but has exact finite relative
norm $-7/288$, rather than zero.  Certificate
\par\noindent
\texttt{REVERSE\_PHYSICS\_BT\_COVARIANT\_GHOST\_}
\allowbreak\texttt{PARITY\_BRANCH\_OBSTRUCTION\_V1}
\par\noindent
therefore preserves the Eq.~(19) charge formula with $Q=0$ while obstructing
the ghost-even neutral term used by its positivity argument on this declared
public branch.  The minimal symmetric generator
$\frac12(Z^{-1}K_++Z\kappa K_+\kappa)$ does yield an exact all-order finite
ghost-even projector, but its added conjugate branch has no public $R_t$ or
source affiliation.  Neither this repair nor the obstruction is a continuum
probability or a gravitational/BRST result.
The same-chart regular source affiliation is itself obstructed.  On the
off-shell commutative local-symbol algebra, zero-jet augmentation sends
$O=\lambda^{-1}Ze^{\lambda\varphi}$ to the Laurent unit
$\lambda^{-1}Z$, while it sends
$Y=Z^{-1}e^{-\lambda\varphi}
[\Box\varphi+\lambda(\partial\varphi)^2]$ to zero.  The first symbol has
explicit inverse $\lambda Z^{-1}e^{-\lambda\varphi}$; the second cannot be a
unit.  Unit preservation therefore forbids any regular unital ghost-parity
automorphism exchanging them, independently of normalization.  Certificate
\par\noindent
\texttt{REVERSE\_PHYSICS\_BT\_PERTURBATIVE\_GHOST\_}
\allowbreak\texttt{PARITY\_UNIT\_OBSTRUCTION\_V1}
\par\noindent
does not impose that augmentation on the CCR/Weyl algebra: singular,
nonlocal or unbounded operator correspondences remain open.  Localization
at the bracketed field strength loses the zero-jet vacuum chart, while a
doubled sheet changes the source theory.  Thus no regular
perturbative-vacuum pullback supplies the missing branch.
The scalar history labels admit a more physical reduced-mode lift.  The
rooted-comb sectors have dimensions $3$, $12$, and $60$, and their exact tree
weights factor into positive per-extension rate squares
$1/48$, $5/64$, and $27/400$.  The certificate
\texttt{REVERSE\_PHYSICS\_BT\_CHANNEL\_}
\allowbreak\texttt{RESOLVED\_BRANCHING\_}
\allowbreak\texttt{INSTRUMENT\_V1}
constructs a finite completely positive trace-preserving branching instrument
whose three first-level jumps have the physical scalar Gram $(1/48)I_2$.  The
higher identity-species lift and three-emission absorbing closure are not
derived scalar amplitude phases or an all-order Hamiltonian.  In fact, the
pre-trace six-point scalar amplitude obstructs that identity lift at the
second jump.  In the constant/linear parent-jet basis its unique coefficients
are $(2Q_{\rm inner},a_2/2)$, rather than
$(2Q_{\rm inner},2L_{\rm inner})$, and the raised spectator-profile
endomorphism has characteristic discriminant
\[
 -\frac{4u^2a_2^2(\tau_2+a_2)(2\tau_2-a_2)}
 {(\tau_2-2a_2)^2}<0
 \qquad (\tau_2>a_2>0).
\]
The certificate
\texttt{REVERSE\_PHYSICS\_BT\_SIX\_POINT\_}
\allowbreak\texttt{PARENT\_JET\_INTERFERENCE\_V1}
therefore retains the scalar $5/3072$ history weight and the abstract finite
instrument but refutes its amplitude affiliation above the first jump on the
two-species carrier.  The grading-faithful parent-jet times spectator-profile
enlargement resolves that scoped obstruction.  On its four-component tensor
carrier the complete pullback obeys
\[
 A^2=2uvA,\qquad P=\frac{A}{2uv}=P^2=P^\sharp .
\]
The kernel of $P$ is nondegenerate and exactly invisible to the physical
collapse, while its orthogonal two-dimensional image has positive scalar
raised Gram $2uvI_2$.  The pointwise identity
\[
 (RDX)^T(3J)(RDX)=2uv(PX)^T(J\otimes3J)(PX)
\]
reconstructs the complete scalar six-point amplitude.  Together with the
five-point prefix replay this affiliates the conditional second rate
$5/64$ on the quotient fibre.  The certificate
\par\noindent
\texttt{REVERSE\_PHYSICS\_BT\_}\allowbreak
\texttt{SIX\_POINT\_}\allowbreak
\texttt{PROFILE\_QUOTIENT\_COMPLETION\_V1}
\par\noindent
is followed by the complete scalar seven-point pre-trace calculation.  Its
seven square-free spectator components reduce to singleton and
complementary-pair profiles and factor uniquely through the same recombined
parent carrier.  The two parent coefficients have opposite signs throughout
the nested physical threshold domain.  After the seven-leg delta-prime sign,
the nonzero physical quotient eigenvalue is $-2uv>0$, and the signed
arbitrary-vector identity reconstructs the complete scalar amplitude before
its final trace.  Consequently
\[
 q_2=\frac{9/81920}{5/3072}=\frac{27}{400}
\]
is amplitude-affiliated rather than merely inserted as a positive identity
lift.  The certificate
\par\noindent
\texttt{REVERSE\_PHYSICS\_BT\_}\allowbreak
\texttt{SEVEN\_POINT\_}\allowbreak
\texttt{PROFILE\_QUOTIENT\_AFFILIATION\_V1}
\par\noindent
checks all $2485$ trees and the signed Krein quotient independently.  This is
still a scalar \texttt{LOCAL-}\allowbreak\texttt{ALGEBRAIC},
\texttt{REDUCED-}\allowbreak\texttt{MODE} result.  Neither
the scalar quotient, the growing history carrier, its three affiliated jump
rates, nor its threshold sign proof transfers to metric BV data; no
gravitational or Lorentzian claim is made.
On that scalar quotient alone, the three jumps further close into a common
finite Krein-skew coupling jet with edge weights
$\sqrt3/12$, $\sqrt{10}/8$, and $9/20$.  Its exponential reproduces the
scalar leading coefficients through three emissions when $x^2=a$, but this
square-root relation proves that it is not generated strongly in additive
resolution on any fixed finite bounded carrier.  Certificate
\texttt{REVERSE\_PHYSICS\_BT\_THREE\_JUMP\_}
\allowbreak\texttt{KREIN\_MOLLER\_JET\_V1}
classifies the finite scalar Taylor jet and the obstruction separately.  The
ordinary generator remains obstructed, but a scalar quantum-stochastic
continuation now exists.  On the same $152$-dimensional quotient, $75$
rooted-comb edge operators drive a bounded Hudson--Parthasarathy cocycle
\cite{HudsonParthasarathy1984}; its vacuum reduction is the exact branching
instrument and its first three ordered-noise sectors reproduce the scalar
M\o ller jet.  The certificate
\par\noindent
\texttt{REVERSE\_PHYSICS\_BT\_QUANTUM\_}
\allowbreak\texttt{STOCHASTIC\_MOLLER\_DILATION\_V1}
\par\noindent
proves minimal edge-noise multiplicity $75$ for the pinned scalar GKSL map
and verifies both unitary structure identities independently.  None of this
carrier, its auxiliary resolution noise, its rates, or its stochastic
cocycle is imported into the metric BV complex.  It therefore establishes
no gravitational or Lorentzian M\o ller statement, and supplies neither a
metric fourth jump nor a spacetime LSZ operator.
The scalar first-emission cocycle now also has an exact Abel-regularized
physical-range intertwiner: polar normalization retains both external-jet
species and the full collinear direct-integral shape, while a noise-only map is
rank-obstructed.  The certificate
\par\noindent
\texttt{REVERSE\_PHYSICS\_BT\_ABEL\_FOCK\_}
\allowbreak\texttt{PHYSICAL\_INTERTWINER\_V1}
\par\noindent
affiliates the first three scalar edge marks.  The six-point positive quotient
also defines an exact cumulative K\"all\'en resolution coordinate and an
ordered two-noise direct-integral isometry for the twelve second-level marks.
The certificate
\par\noindent
\texttt{REVERSE\_PHYSICS\_BT\_SIX\_POINT\_NESTED\_}
\allowbreak\texttt{CONTINUUM\_INTERTWINER\_V1}
\par\noindent
extends scalar physical continuum affiliation to marks $0$ through $14$.
At seven points the signed quotient eigenvalue times the massless-hierarchy
middle K\"all\'en measure defines a second exact cumulative coordinate with an
elementary primitive and unit linear asymptote.  The certificate
\par\noindent
\texttt{REVERSE\_PHYSICS\_BT\_SEVEN\_POINT\_NESTED\_}
\allowbreak\texttt{CONTINUUM\_INTERTWINER\_V1}
\par\noindent
composes its conditional column with the two-emission range and physically
affiliates marks $15$ through $74$.  All 75 marks in the available scalar
three-emission hierarchy therefore have continuum columns, but no fourth jump
or all-order inductive operator is constructed.  Composing their direct sum
with the scalar 75-channel HP unitary and vacuum inclusion gives an exact
finite physical continuum column
${\cal M}_a={\cal A}_{\leq3}U_aI_\Omega$ with
${\cal M}_a^*{\cal M}_a=I_2$.  Its nonnegative hard and three emitted-sector
norms sum to one for every resolution length.  In that pinned finite model
the hard response $-1/16$ and real response $+1/16$ therefore cancel on the
same physical column; this does not derive the full scalar NLO dynamics.
The accompanying first-emission compression theorem retains the public
rank-one nilpotent $R_t$ leg only after a forced two-dimensional cross-Krein
complement whose covariant determinant is $-\rho^2$ and whose raised trace
is $-2\rho$.  A positive-noise or trace-null-only completion is impossible.
Certificate
\par\noindent
\texttt{REVERSE\_PHYSICS\_BT\_FINITE\_PHYSICAL\_}
\allowbreak\texttt{MOLLER\_COLUMN\_V1}
\par\noindent
keeps that compression distinct from equality with the physical splitting
map and does not derive the missing block from BT dynamics.  The fixed
endpoint freedom does not supply the block universally on its minimal scalar
two-profile lift.  Relative to the symmetric triple-plus reference, exact
matching forces
$c_0=0$, $c_1=7/4+\rho/2$, and $c_2=-7/8$; three physical rational fixtures
have distinct $\rho$ and therefore require distinct $c_1$.  This
reference-independent slope obstruction is certified by
\par\noindent
\texttt{REVERSE\_PHYSICS\_BT\_ENDPOINT\_COMPLEMENT\_}
\allowbreak\texttt{MATCHING\_V1}.
\par\noindent
It rules out one fixed scalar endpoint prescription, not
spectator-dependent scalar dynamics, and none of its probe lift or matching
coefficients transfers to the metric BV carrier.  Neither the
scalar polar ranges, the cumulative resolution measures, the vacuum column,
nor the minimal cross-Krein complement transfers to the metric BV complex.
Even in the scalar problem, the vacuum column does not determine a scattering
operator on arbitrary incoming states.  The exact defect-completion theorem
parameterizes every same-space unitary extension as
$S_W={\cal M}_aI^*+W(1-II^*)$, where $W$ is a partial unitary from the incoming
defect onto $1-{\cal M}_a{\cal M}_a^*$.  Two different finite exact
extensions have the same column, and the physical continuum defect is
countably infinite.  The certificate
\texttt{REVERSE\_PHYSICS\_BT\_PHYSICAL\_MOLLER\_}
\allowbreak\texttt{DEFECT\_COMPLETION\_V1}
thus proves abstract scalar unitary extendibility and, simultaneously, exact
underdetermination of the incoming-continuum dynamics.  It supplies no metric
defect space, gravitational asymptotic Hamiltonian, LSZ map, or causal claim.
There is nevertheless a genuine scalar hard-scattering baseline.  The
certificate
\texttt{REVERSE\_PHYSICS\_BT\_UV\_HARD\_}\allowbreak
\texttt{SCATTERING\_LAW\_V1}
proves at leading logarithmic order that every fixed nonforward scalar
acceptance falls as $24\pi^2/[25s\log^2(s/\Lambda^2)]$ per unit solid angle.
This uses the scalar Born normalization, scalar beta function, and scalar
projected hard log.  None of its coefficient, positivity statement, or
Callan--Symanzik closure transfers to the graviton/BRST carrier; the gravity
paper records it only as the kind of physical endpoint an independent
gravitational construction would have to reach.

\section{Discussion}\label{sec:discussion}

\subsection{The same order pattern, for different reasons}

Paper~V and the present gravity calculation share the sequence
\begin{equation}
\text{first-order protection}\longrightarrow
\text{second-order scattering obstruction}.
\end{equation}
The agreement in order should not obscure the difference in structure.
The perfect-square scalar is protected by its special cubic kinematic
factor and reality assignments.  Gravity is protected because the
Einstein solution space is an exact nonlinear subsector and the remaining
odd cubic vertex is kinematically closed.  In the scalar, an enhanced
$O(1,1)$ charge appears exactly at the perfect-square boundary and can
sort an obstruction into a one-sided null ideal.  In gravity, the $MMM$
and $MMh$ interactions remove every uniform abelian charge before the
four-point obstruction is considered.

\subsection{Relation to the PT-symmetric quadratic-gravity proposal}

Kuntz \cite{KuntzPTGravity} starts from the same scalar-free
Einstein--Weyl vertex content used here.  His auxiliary-field complex
deformation makes the massless and massive free spin-$2$ sectors
Hermitian and positive and places factors of $i$ on interactions with an
odd number of massive fields.  At this level his rotation and
\eqref{eq:quarterturn} are the same construction up to the sign chosen
for the quarter-turn.  His transverse physical asymptotic space also
retains both massless helicities and all five massive polarizations, so
the distinction is not that the $M$ states have been removed.

The difference is in the interacting test.  The displayed metric
calculation in Ref.~\cite{KuntzPTGravity} replaces the full tensor field
theory by a mode Hamiltonian with spin and momentum dependence suppressed
and with schematic interaction coefficients.  It demonstrates a
positive perturbative metric for that truncation and offers it as
evidence for the gravitational theory, but it does not include the exact
energy-degenerate $MM\to Mh$ block.  Proposition~\ref{prop:bornidentity}
shows why that omission is decisive: a nonzero anti-Hermitian component
inside a degenerate free-energy block is not divided by a free-energy
difference and cannot be removed by an analytic $R_2$.

Thus our result is a direct full-vertex stress test of that proposal, not
a general rejection of PT-symmetric gravity.  Any extension of the
truncated construction to the complete theory must evade
\eqref{eq:obstruction} by leaving at least one of the present hypotheses:
for example through a nonanalytic metric, a different free base point,
different asymptotic states, or a prescription that excludes the massive
mode from the physical scattering spectrum.

\subsection{Real forms become dynamical}

At the free level a real form selects which complex solutions are called
real and which adjoint defines observables.  The choice can appear
kinematic.  The phase theorem makes its interacting content explicit:
every physical process with odd external massive number crosses the two
reality assignments and becomes anti-Hermitian in the positive frame.
The decisive diagnostic is therefore
\begin{equation}
\boxed{\text{Which physical on-shell processes change the real-form
sector?}}
\end{equation}
For Einstein--Weyl gravity the first answer is $MM\to Mh$.

\subsection{Why the factorization and physical rails are both useful}

The Einstein-frame residue \eqref{eq:residue} is compact, independent,
and immune to cancellation by local contact terms.  It proves a genuine
interaction before the large quartic calculation is attempted.  The
original-variable physical point then supplies what factorization alone
does not: an exact real-shell value, the full gauge cancellation, the
physical adjoint, and the deformation-complex matrix element.  Agreement
between the two rails is substantially stronger than either calculation
in isolation.

\subsection{What Krein visibility means}

The negative traces in \eqref{eq:Xtrace} and \eqref{eq:Osharp} are not
advertised as negative probabilities.  They show that the conventional
quadratic Krein form sees the block: it is neither zero nor killed by a
grading selection rule.  A viable interacting probability construction
would need extra structure beyond the free fundamental symmetry.  The
scalar boundary supplies a continuous one-sided charge at tree level, but the
exact scalar loop calculation shows that this algebraic structure alone does
not select an ordinary axis-compatible infrared prescription.  The finite scalar
projector transport shows that coherent cross-multiplicity
normalization is algebraically compatible after Born-rate normalization,
while the ordinary cubic Fock generator has zero collinear Gram; it does not
by itself show that the required Jordan term exists.  The exact scalar
calculation derives its two-annihilator coefficient and proves secular
cancellation, but stops at a cubic endpoint-pole obstruction before a
continuum projector or probability exists.
The verified
Einstein--Weyl vertices do not even supply the corresponding uniform charge,
and no gravitational infrared completion is constructed here.

The unresolved gravitational problem is whether conformal-boundary conditions,
changed observable algebras, or genuinely nonlocal gravitational
structures can evade the hypotheses of Theorem~\ref{thm:master} without
destroying covariance and factorization.  Any proposal that retains
massless asymptotic gravitons must additionally declare how its real and
virtual collinear sectors are glued and prove regulator independence on the
BRST quotient; the scalar certificate motivates this requirement but is not
evidence that gravity passes or fails it.

\begin{thebibliography}{99}

\bibitem{Stelle}
K.~S.~Stelle,
\emph{Renormalization of higher-derivative quantum gravity},
Phys.\ Rev.\ D \textbf{16} (1977) 953,
\href{https://doi.org/10.1103/physrevd.16.953}{doi:10.1103/physrevd.16.953}.

\bibitem{Salvio}
A.~Salvio,
\emph{Quadratic gravity},
Front.\ Phys.\ \textbf{6} (2018) 77,
arXiv:1804.09944.

\bibitem{HindawiOvrutWaldram}
A.~Hindawi, B.~A.~Ovrut and D.~Waldram,
\emph{Consistent spin-two coupling and quadratic gravitation},
Phys.\ Rev.\ D \textbf{53} (1996) 5583--5596,
arXiv:hep-th/9509142.

\bibitem{MagnanoSokolowski}
G.~Magnano and L.~M.~Soko\l{}owski,
\emph{Nonlinear massive spin-two field generated by higher derivative
gravity},
Ann.\ Phys.\ \textbf{306} (2003) 1--35,
arXiv:gr-qc/0209022.

\bibitem{DonaEtAl}
P.~Don\`a, S.~Giaccari, L.~Modesto, L.~Rachwa\l{} and Y.~Zhu,
\emph{Scattering amplitudes in super-renormalizable gravity},
JHEP \textbf{08} (2015) 038,
arXiv:1506.04589.

\bibitem{MostafazadehMetric}
A.~Mostafazadeh,
\emph{Metric operator in pseudo-Hermitian quantum mechanics and the
imaginary cubic potential},
J.\ Phys.\ A \textbf{39} (2006) 10171--10188,
arXiv:quant-ph/0508195.

\bibitem{BBJQFT}
C.~M.~Bender, D.~C.~Brody and H.~F.~Jones,
\emph{Extension of PT-symmetric quantum mechanics to quantum field
theory with cubic interaction},
Phys.\ Rev.\ D \textbf{70} (2004) 025001,
arXiv:hep-th/0402183.

\bibitem{KuboKuntz}
J.~Kubo and J.~Kuntz,
\emph{Analysis of unitarity in conformal quantum gravity},
Class.\ Quantum Grav.\ \textbf{39} (2022) 175010,
arXiv:2202.08298.

\bibitem{KuntzPTGravity}
J.~Kuntz,
\emph{Unitarity through PT symmetry in quantum quadratic gravity},
arXiv:2410.08278.

\bibitem{AdamoNakachTseytlin}
T.~Adamo, S.~Nakach and A.~A.~Tseytlin,
\emph{Scattering of conformal higher spin fields},
JHEP \textbf{07} (2018) 016,
arXiv:1805.00394.

\bibitem{JohanssonMogullTeng}
H.~Johansson, G.~Mogull and F.~Teng,
\emph{Unraveling conformal gravity amplitudes},
JHEP \textbf{09} (2018) 080,
arXiv:1806.05124.

\bibitem{DonoghueMenezes}
J.~F.~Donoghue and G.~Menezes,
\emph{Unitarity, stability and loops of unstable ghosts},
Phys.\ Rev.\ D \textbf{100} (2019) 105006,
arXiv:1908.02416.

\bibitem{AnselmiFakeon}
D.~Anselmi,
\emph{Fakeons and Lee--Wick models},
JHEP \textbf{02} (2018) 141,
arXiv:1801.00915.

\bibitem{LiuLuPope}
H.-S.~Liu, H.~L\"u, C.~N.~Pope and J.~F.~V\'azquez-Poritz,
\emph{Not conformally-Einstein metrics in conformal gravity},
arXiv:1303.5781.

\bibitem{KOS}
S.~Kamefuchi, L.~O'Raifeartaigh and A.~Salam,
\emph{Change of variables and equivalence theorems in quantum field
theories},
Nucl.\ Phys.\ \textbf{28} (1961) 529--549,
\href{https://doi.org/10.1016/0029-5582(61)90056-6}{doi:10.1016/0029-5582(61)90056-6}.

\bibitem{WeinbergQFT}
S.~Weinberg,
\emph{The Quantum Theory of Fields, Vol.~I: Foundations},
Cambridge University Press, 1995.

\bibitem{BM2008}
C.~M.~Bender and P.~D.~Mannheim,
\emph{No-ghost theorem for the fourth-order derivative Pais--Uhlenbeck
oscillator model},
Phys.\ Rev.\ Lett.\ \textbf{100} (2008) 110402,
\href{https://doi.org/10.1103/physrevlett.100.110402}{doi:10.1103/physrevlett.100.110402}.

\bibitem{BM2008PRD}
C.~M.~Bender and P.~D.~Mannheim,
\emph{Exactly solvable PT-symmetric Hamiltonian having no Hermitian
counterpart},
Phys.\ Rev.\ D \textbf{78} (2008) 025022, arXiv:0804.4190.

\bibitem{BT2026}
S.~Bateman and N.~Turok,
\emph{Escape from Ostrogradsky via hidden ghost parity},
arXiv:2607.00096.

\bibitem{HudsonParthasarathy1984}
R.~L.~Hudson and K.~R.~Parthasarathy,
\emph{Quantum It\^o's formula and stochastic evolutions},
Commun.\ Math.\ Phys.\ \textbf{93} (1984) 301--323,
\href{https://doi.org/10.1007/BF01258530}{doi:10.1007/BF01258530}.

\bibitem{Maldacena}
J.~Maldacena,
\emph{Einstein gravity from conformal gravity},
arXiv:1105.5632.

\bibitem{AzizovIokhvidov}
T.~Ya.~Azizov and I.~S.~Iokhvidov,
\emph{Linear Operators in Spaces with an Indefinite Metric},
Wiley, 1989.

\bibitem{paperI}
GPT-5.6.sol and A.~A.~Palm,
\emph{Canonical Positive Symplectic Diagonalization of the
Pais--Uhlenbeck Oscillator},
Paper~I of the pure-Weyl gravity programme, 2026;
file \path{paper/01-symplectic-diagonalization.pdf} in the repository
\url{https://github.com/area9innovation/weyl-gravity}.

\bibitem{paperII}
GPT-5.6.sol and A.~A.~Palm,
\emph{The Pais--Uhlenbeck Metric as a Minimum-Distortion Principle,
and the Representation Problem for the Fourth-Order Field},
Paper~II of the pure-Weyl gravity programme, 2026;
file \path{paper/02-variational-fock.pdf} in the repository
\url{https://github.com/area9innovation/weyl-gravity}.

\bibitem{paperIII}
GPT-5.6.sol and A.~A.~Palm,
\emph{The Universal Vacuum of the Fourth-Order Scalar Field:
Metric Orbits, Fock Sectors, and the Krein Boundary},
Paper~III of the pure-Weyl gravity programme, 2026;
file \path{paper/03-fourth-order-vacuum.pdf} in the repository
\url{https://github.com/area9innovation/weyl-gravity}.

\bibitem{paperIV}
GPT-5.6.sol and A.~A.~Palm,
\emph{Gauge Reduction and the Completion Problem in Fourth-Order Gravity:
PU Pairing, Covariant Real Forms, and the Conformal Jordan Boundary},
Paper~IV of the pure-Weyl gravity programme, 2026;
file \path{paper/04-fourth-order-gravity.pdf} in the repository
\url{https://github.com/area9innovation/weyl-gravity}.

\bibitem{paperV}
GPT-5.6.sol and A.~A.~Palm,
\emph{Interaction Obstructions, Resonant PT Breaking,
and Doubled Jordan Symmetry in Fourth-Order Theories},
Paper~V of the pure-Weyl gravity programme, 2026;
file \path{paper/05-interaction-obstructions.pdf} in the repository
\url{https://github.com/area9innovation/weyl-gravity}.

\end{thebibliography}

\appendix

\section{Einstein consistent truncation and LSZ}
\label{app:truncation}

This appendix spells out the hypotheses behind
Lemma~\ref{lem:onemassive}.  For the scalar-free combination in
\eqref{eq:action}, the field equation is the Einstein tensor plus a
multiple of the Bach tensor.  In four dimensions the Bach tensor may be
written
\begin{equation}
B_{\mu\nu}=\left(\nabla^\rho\nabla^\sigma
+\frac12R^{\rho\sigma}\right)C_{\mu\rho\nu\sigma}
\end{equation}
up to a convention-dependent overall sign.  It vanishes on every
Einstein metric \cite{LiuLuPope}.  In the flat-background problem the
relevant exact subsector is Ricci-flat vacuum Einstein gravity.

On the regular branch of the Legendre transformation from Ricci-tensor
gravity to its second-order form, introduce the Einstein metric and a
symmetric massive field $\Phi_{\mu\nu}$.  The transformation is algebraic
in the auxiliary variables and locally invertible wherever its determinant
does not vanish \cite{MagnanoSokolowski}.  Linearization about the common
flat vacuum diagonalizes the asymptotic equations into
\begin{equation}
\mathcal K_hh=0,
\qquad
\mathcal K_MM=0,
\end{equation}
with $p^2=0$ and $p^2=M^2$, respectively.  Thus $h$ and $M$ used in the
main text are LSZ mass eigenfields, not arbitrary components of the
original metric perturbation.

Let $E_M[h,M]=0$ be the exact massive-field equation in this frame.  The
consistent truncation states
\begin{equation}\label{eq:truncapp}
E_M[h_E,0]=0
\end{equation}
for every nonlinear Einstein solution $h_E$.  Expanding about flat space,
\begin{equation}
E_M[h,M]=\mathcal K_MM+
\sum_{n\ge2}J_M^{(n)}[h^n]
+\sum_{r\ge1,s\ge1}J_{r,s}[M^rh^s].
\end{equation}
Equation \eqref{eq:truncapp} does not require each off-shell functional
$J_M^{(n)}$ to vanish.  It requires its value on a full perturbative
Einstein solution to cancel, including nonlinear corrections to that
solution.  LSZ amputation removes precisely the external inverse-
propagator and field-redefinition terms, leaving
\begin{equation}
\mathcal A(M,h_1,\ldots,h_n)=0
\end{equation}
for physical on-shell Einstein gravitons.  This is the standard
consistent-truncation selection rule.

The perturbiner supplies an explicit cubic check.  In signature
$(+,-,-,-)$, take all momenta incoming,
\begin{equation}
p_M=(1,0,0,0),\qquad
p_2=\left(-\frac12,0,0,-\frac12\right),\qquad
p_3=\left(-\frac12,0,0,\frac12\right).
\end{equation}
They obey $p_M^2=1$, $p_2^2=p_3^2=0$ and sum to zero.  The exact action
coefficient vanishes for all five massive TT polarizations and the four
products of the two graviton helicities.  Replacing either graviton
polarization by $p_{(\mu}\xi_{\nu)}$ also gives zero.

The selection rule is tree-level and concerns external mass eigenfields.
It does not state that the original fourth-order metric action lacks
terms apparently linear in a chosen massive component, nor does it
exclude loop effects, anomalies, or a singular failure of the
mass-eigenfield map at the conformal boundary.

\section{Einstein-frame cubic expansion}
\label{app:eframe}

We record enough of the independent G14 calculation to make the exact
numbers reproducible.  This appendix uses signature $(-,+,+,+)$ and
$M=1$.  It matches
\path{symbolic/verify_gravity_factorization.py}.

For a symmetric matrix $\Phi$, write
\begin{equation}
X=g^{-1}\Phi,\qquad
t_n=\operatorname{tr}(X^n),\qquad
A=\sqrt{\det(I+X)}.
\end{equation}
Through cubic order,
\begin{align}
A={}&1+\frac12t_1+\frac18t_1^2-\frac14t_2
+\frac1{48}t_1^3-\frac18t_1t_2+\frac16t_3+O(\Phi^4),\nonumber\\
A^{-1}={}&1-(A-1)+(A-1)^2-(A-1)^3+O(\Phi^4).
\end{align}
In the potential of \eqref{eq:LM}, $\tau=4+t_1$ and the relevant
quadratic contraction is $4+2t_1+t_2$.  Substitution gives, when
$t_1=0$,
\begin{equation}
A^{-1}\left(-\tau^2+\psi_{\mu\nu}\psi^{\mu\nu}
+6A\tau-12A^2\right)=t_2+O(\Phi^4),
\end{equation}
with zero $t_3$ coefficient.  This is the potential cancellation
\eqref{eq:potentialcancel}.

At first order in one plane wave $(p,E)$,
\begin{equation}
Q^{(1)\alpha}{}_{\mu\nu}(E,p)
=\frac{i}{2}\eta^{\alpha\beta}
\left(p_\mu E_{\beta\nu}+p_\nu E_{\beta\mu}
-p_\beta E_{\mu\nu}\right).
\end{equation}
The quadratic correction $Q^{(2)}$ comes from the inverse-$\psi$
factor.  The cubic $MMM$ coefficient is the sum of all contractions
$K(Q^{(1)},Q^{(2)})+K(Q^{(2)},Q^{(1)})$, with the factor $1/2$ in
\eqref{eq:LM}.  Exact substitution of \eqref{eq:mmmkin} and normalized
TT tensors gives \eqref{eq:mmmvalue}.

For $MMh$, the script expands $g^{-1}$, $\psi^{-1}$, the connection of
$g$, and $Q$ through the orders required for the coefficient of three
independent wave labels.  At \eqref{eq:mmhkin}, take
\begin{equation}
v_r=(1,-1,0,0),\qquad e_h=(1,0,0,1),
\qquad E_r=v_r^\flat\otimes v_r^\flat,
\qquad H=e_h^\flat\otimes e_h^\flat.
\end{equation}
Together with the internal $+$ polarization at rest these tensors are
transverse and traceless, and the coefficient is \eqref{eq:mmhvalue}.
For
\begin{equation}
H^{\rm gauge}_{\mu\nu}=k_\mu\xi_\nu+k_\nu\xi_\mu
\end{equation}
with four independent symbolic components of $\xi$, the coefficient
vanishes identically.

Finally, the massive TT two-point coefficient for waves $(p,E)$ and
$(-p,E)$ with $E\cdot E=1$ is
\begin{equation}
L_M^{(2)}=\frac{p^2+1}{4}.
\end{equation}
This fixes the factor of four between the polarization numerator
$\sqrt3/32$ and pole residue $\sqrt3/8$; it is not inferred from a
unit-residue completeness convention.

\section{Polarization basis and factorization sum}
\label{app:polarizations}

At the rest momentum $P=(1,0,0,0)$ in the Einstein-frame signature,
let $e_x,e_y,e_z$ be unit spatial vectors.  A normalized massive basis is
\begin{align}
E_+&=\frac1{\sqrt2}(e_xe_x-e_ye_y),
&E_\times&=\frac1{\sqrt2}(e_xe_y+e_ye_x),\nonumber\\
E_0&=\frac1{\sqrt6}(e_xe_x+e_ye_y-2e_ze_z),
&E_{xz}&=\frac1{\sqrt2}(e_xe_z+e_ze_x),\nonumber\\
&&E_{yz}&=\frac1{\sqrt2}(e_ye_z+e_ze_y).
\end{align}
Each is transverse, traceless, and orthonormal under the induced tensor
inner product.  For a massive momentum $(E,0,0,p_z)$, replace $e_z$ by
the normalized longitudinal transverse vector; the same formulas give
the transported helicity basis.

The factorization numerator is basis-independent.  In a nonorthonormal
basis $T_a$ with Gram matrix $G_{ab}=T_a\cdot T_b$, it is
\begin{equation}
N=V_a(G^{-1})^{ab}U_b,
\end{equation}
where $V_a=\mathcal A_3(M,M,T_a)$ and
$U_b=\mathcal A_3(T_b,M,h)$.  The orthonormal basis above reduces this
to the sum displayed after \eqref{eq:resnum}.  All five intermediate
states are included; the result is not one selected nonzero channel.

The original fourth-order perturbiner supplies a separate complex
factorization check at a rational-$i$ point.  It constructs an exact
five-dimensional TT basis at the internal massive momentum, contracts
the two cubic vectors with its Gram inverse, and obtains a nonzero
reduced numerator.  The Einstein-frame calculation then fixes a compact
normalized representative and the explicit inverse-kernel slope.

\section{Physical-shell geometry and the rational point}
\label{app:shell}

For incoming equal masses in the center-of-mass frame,
\begin{equation}
p_1=(E,0,0,p),\qquad p_2=(E,0,0,-p),
\qquad E=\frac{\sqrt s}{2},\quad
p=\frac12\sqrt{s-4M^2}.
\end{equation}
Choose the outgoing graviton at scattering angle $\theta$,
\begin{equation}
k=(\kappa,\kappa\sin\theta,0,\kappa\cos\theta),
\qquad
\kappa=\frac{s-M^2}{2\sqrt s},
\end{equation}
and $p_3=p_1+p_2-k$.  Then $p_3^2=M^2$ and
$p_3^0=(s+M^2)/(2\sqrt s)$.  The domain
\begin{equation}
s>4M^2,\qquad -1<\cos\theta<1
\end{equation}
is a real open chart away from forward/backward and threshold
degeneracies.

At $s=25M^2/4$ one has $E=5M/4$ and
$\kappa=21M/20$.  Taking $\cos\theta=3/5$ and
$\sin\theta=4/5$ gives exactly \eqref{eq:rationalpoint}.  All energy,
momentum, and mass-shell identities are rational.  Because all spatial
$y$ components vanish, reflection $y\mapsto-y$ is a symmetry.  An odd
number of $y$-odd polarization tensors forces an amplitude zero.  The
certificate uses the first nonzero all-even combination.

The exact threshold regression uses two incoming massive particles at
rest and an outgoing graviton of energy $3M/4$.  In the separate
polarization convention of that test,
\begin{equation}
\mathcal A_{\rm threshold}=-\frac{509784}{390625}.
\end{equation}
Threshold is not used for the open-set claim; it guards the implementation
against drift.

\section{Full four-point engine and deformation certificate}
\label{app:fourpoint}

\subsection{Wave algebra}

A wave-algebra element is a dictionary from subsets of external labels
to sparse tensor components.  Multiplication discards products in which
a wave label repeats, so multilinear truncation is automatic.  A
derivative acting on subset $S$ multiplies by
\begin{equation}
iP_{S\mu}=i\sum_{a\in S}p_{a\mu}.
\end{equation}
From the metric, inverse metric, determinant, Christoffel symbols, and
Ricci tensor, the engine extracts
\begin{equation}
\sqrt{-g}R,\qquad
\sqrt{-g}R_{\mu\nu}R^{\mu\nu},
\qquad
\sqrt{-g}R^2
\end{equation}
at the desired wave subset.  Before any cubic or quartic claim it checks
that the two-point coefficient vanishes on $p^2=0$ and $p^2=M^2$, is
nonzero off shell, and obeys cubic Ward identities.

\subsection{Exact polarization tensors}

For a massive momentum $p$ in the $(t,x,z)$ plane, choose two independent
transverse in-plane vectors $v_1,v_2$.  With
\begin{equation}
\Pi_{\mu\nu}=\eta_{\mu\nu}-\frac{p_\mu p_\nu}{M^2},
\end{equation}
the parity-even TT tensor used in the certificate is
\begin{equation}
E_{12,\mu\nu}=v_{1\mu}v_{2\nu}+v_{2\mu}v_{1\nu}
-\frac23(v_1\cdot v_2)\Pi_{\mu\nu}.
\end{equation}
The three massive external legs use the corresponding $E_{12}$ at
$p_1,p_2,-p_3$.  For the graviton momentum direction proportional to
$(5,4,0,3)$, use transverse vectors
\begin{equation}
e_1=(0,-3,0,4),\qquad e_2=(0,0,1,0),
\end{equation}
and the real plus tensor
\begin{equation}
\varepsilon_+=\frac1{25}e_1e_1-e_2e_2.
\end{equation}
These are the convention-fixed polarizations entering \eqref{eq:AK}.

\subsection{Gauge and adjoint checks}

For each exchange, the exact bordered solution obeys
\begin{equation}
CX=0,
\qquad
KX+C^T\lambda+J=0.
\end{equation}
The total amplitude with graviton gauge polarization is zero, while the
contact contribution alone is
\begin{equation}
-\frac{596335802837691}{361367187500000}\ne0.
\end{equation}
This cancellation is a sensitive check of relative signs and
combinatorics.  De Donder and axial constraints give identical physical
totals.

Under reversal $p_a\mapsto-p_a$, all contacts and cubic currents are
recomputed.  The action contains two- and four-derivative terms, and the
independently recomputed quadratic matrices obey $K(-P)=K(P)$.  Contact,
$s$, $t$, and $u$ pieces equal their forward values separately.  Since
the external linear tensors are real, this is the physical reverse
matrix element used in $\Born_{2,+}^{\dg}$.

\subsection{Pole regression}

An exact shell-preserving complex shift
\begin{equation}
p_1(z)=p_1+z\eta,\qquad k(z)=k-z\eta,
\qquad \eta=(3,0,4i,5),
\end{equation}
has $\eta^2=\eta\cdot p_1=\eta\cdot k=0$.  In the relevant channel
\begin{equation}
P(z)^2=\frac{25}{4}+15z
\end{equation}
hits the massive pole at $z_*=-7/20$.  The bordered kernel there has
exactly the five TT zero modes.  If $T$ collects them and $G_5$ is their
Gram matrix,
\begin{equation}
W=T^TK'(z_*)T=-\frac{15}{2}G_5
\end{equation}
entrywise, including 24 nontrivial identities.  The residue
\begin{equation}
-j_1^TW^{-1}j_2=-\frac{4551434384\,i}{328515625}\ne0
\end{equation}
equals the independently normalized cubic contraction divided by the
kinetic slope.  A direct near-pole numerical probe of the actual bordered
exchange agrees to relative error below $10^{-4}$.  This regression is
not used to set the convention of the compact Einstein-frame residue;
it validates the original-variable propagator machinery.

\section{Covariant grading and Fock-lift equations}
\label{app:grading}

The exact grading check solves the commutant rather than assuming
Schur's lemma numerically.
At a standard massive momentum, represent the spin-$2$ fiber by symmetric
traceless $3\times3$ tensors.  The three $SO(3)$ generators act as
\begin{equation}
J_i(T)=L_iT+TL_i^T.
\end{equation}
Solving $[X,J_i]=0$ for a general $5\times5$ matrix gives
$X=a_MI_5$.  On the two complex massless helicity fibers, the little-
group rotation generator is $\operatorname{diag}(2,-2)$, whose
commutant is diagonal.  Solving simultaneously with the mass Casimir on
the full seven-dimensional fiber gives
\begin{equation}
X=\operatorname{diag}(a_+,a_-,a_MI_5).
\end{equation}
Hermiticity and $X^2=I$ restrict the three scalars to signs.  The
helicity-exchange real structure sets $a_+=a_-$; the free signature fixes
$(a_+,a_-,a_M)=(1,1,-1)$.

For the Fock lift, cluster factorization is the monoid equation
\eqref{eq:monoid}.  Iterating it gives
\begin{equation}
j(n_h,n_M)=j(1,0)^{n_h}j(0,1)^{n_M}=(-1)^{n_M}.
\end{equation}
The exact script verifies multiplicativity for all partitions in a
finite table and the analytic monoid argument establishes it for all
particle numbers.

\section{BRST block and non-nullity}
\label{app:brst}

The finite verification sets the common real external normalization
$\eta_{\rm conv}\mathcal N_{fi}$ to one and therefore uses $\AKred$ in
place of $\AK$.  Restoring the nonzero multiplier rescales every
quadratic trace by its square and cannot change nullity.

The minimal non-null calculation uses
\begin{equation}
J_{\min}=\begin{pmatrix}1&0\\0&-1\end{pmatrix},
\qquad
X=\begin{pmatrix}0&0\\i\AK&0\end{pmatrix}.
\end{equation}
Then
\begin{equation}
X^\sharp=J_{\min}X^{\dg}J_{\min}
=\begin{pmatrix}0&i\AK\\0&0\end{pmatrix},
\qquad
X^\sharp X=\begin{pmatrix}-\AK^2&0\\0&0\end{pmatrix}.
\end{equation}
For $\mathcal O=-2i\AK\sigma_x$,
\begin{equation}
J_{\min}\mathcal O^{\dg}J_{\min}=\mathcal O,
\qquad \mathcal O^2=-4\AK^2I.
\end{equation}

To encode the cohomology argument, take a canonical finite splitting into
two physical representatives and one contractible doublet $u\mapsto
v\mapsto0$.  In the ordered basis $(i,f,u,v)$,
\begin{equation}
Q=\begin{pmatrix}
0&0&0&0\\
0&0&0&0\\
0&0&0&0\\
0&0&1&0
\end{pmatrix},
\qquad Q^2=0.
\end{equation}
For a completely general $4\times4$ matrix $Y$, direct multiplication
shows
\begin{equation}
\bigl(QY+YQ\bigr)_{\rm phys,phys}=0.
\end{equation}
The upper-left physical block of the exact obstruction is nonzero after
embedding, so it cannot be of this exact form.  This finite matrix is not
a replacement for gravitational BRST quantization; it records the
universal algebraic implication used after the Ward and reduced-helicity
checks have established that the amplitude descends to cohomology.

\section{Machine-verification index}
\label{app:checks}

All scripts named below are in \path{symbolic/} and print
\texttt{ALL PASS}.  Checks use SymPy exact arithmetic except for the
explicitly labelled near-pole regression.

\paragraph{Archive.}
The engine, polarization data and manuscript sources are archived in a
GitHub repository:
\begin{center}
\url{https://github.com/area9innovation/weyl-gravity}.
\end{center}
The historical pre-extraction verification snapshot was
\begin{center}
\texttt{aeb9acaa944f53cbf29fba37476bfe888c6a1ddd}
\end{center}
in the former monorepo; current paths are relative to the standalone
repository root.

\paragraph{Environment.}
The recorded versions are Python~3.14.4 and SymPy~1.14.0.  The pinned
dependency and the SHA-256 manifest are, respectively,
\begin{center}
\path{symbolic/paper6-requirements.txt}\\
\path{symbolic/paper6-verification.sha256}.
\end{center}
Momentum, signature, polarization, action and reduced-amplitude
conventions are recorded in
Appendices~\ref{app:eframe}--\ref{app:fourpoint}; they are explicit
script inputs rather than hidden generated defaults.

\begin{center}
\footnotesize
\begin{tabular}{p{0.11\linewidth}p{0.50\linewidth}p{0.29\linewidth}}
\hline
Check & Claim & Artifact/type\\
\hline
G13a--b & two-point on-shell validation and cubic Ward identities &
\path{gravity_perturbiner.py}, exact\\
G13 & $Mhh=0$ for $5\times4$ physical polarization choices &
\path{verify_gravity_cubic.py}, exact\\
G14a--c & nonzero original-variable $MMM$, $MMh$, and reduced massive-pole
factorization & \path{verify_gravity_cubic.py}, exact\\
G14d--i & potential cancellation, compact $-\sqrt6/8$ and $-\sqrt2/8$
amplitudes, symbolic Ward identity, five-state sum, inverse kernel &
\path{verify_gravity_factorization.py}, exact\\
G15a--f & full physical contact+$s,t,u$, Ward, Bose, gauge, and threshold
checks & \path{verify_gravity_g15.py}, exact\\
G15g & shifted-pole kernel, Schur slope, residue; near-pole probe &
same script, exact plus numerical regression\\
G16 & crossing needed for the adjoint is included in G17; exhaustive
$250$-polarization scan is deferred & not a theorem dependency\\
G17a--d & exposed contact/exchanges, reverse process, EOM/Ward/gauge,
termwise quarter-turn & \path{verify_gravity_obstruction.py}, exact\\
G17e--g & exact shell obstruction, first-order metric-ambiguity
independence, and Born--deformation identity & same script, exact\\
G18a--b & one-particle commutant and cluster-factorizing Fock lift &
\path{verify_gravity_krein.py}, exact\\
G18c--e & non-null quadratic traces, charge no-go, and BRST block &
same script, exact\\
\hline
\end{tabular}
\end{center}

The analytic inputs not reducible to finite symbolic identities are
identified in the main text: the induced-representation interpretation
of the free commutant (Paper~IV), the consistent-truncation/LSZ argument,
covariant--stationary tree matching, real-analytic continuation on the
physical shell, and the standard BRST decoupling statement.  The symbolic
checks test every convention-sensitive coefficient consumed by those
arguments.

\paragraph{Reproduction commands.}
From the repository root for this project:
\begin{verbatim}
python3 symbolic/verify_gravity_paper6.py
\end{verbatim}
The expected terminal line is
\texttt{PAPER VI GRAVITY SUITE: ALL PASS}.  The optional
\texttt{--quick} flag skips the slower shifted-pole regression; the
theorem-critical obstruction script independently recomputes the exact
physical certificate.  Individual script outputs provide the canonical
coefficient-level verification evidence mapped in the table above.

\end{document}
